\documentclass[11pt]{article}
% =======================================================================
%   Shared style for CS 300 notes.
%   Usage: place `% =======================================================================
%   Shared style for CS 300 notes.
%   Usage: place `% =======================================================================
%   Shared style for CS 300 notes.
%   Usage: place `\input{../../notes-style.tex}` right after \documentclass
%   in each chapter's lectures.tex and notes.tex.
% =======================================================================

% ---- Layout -----------------------------------------------------------
\usepackage[margin=1in]{geometry}
\usepackage{parskip}        % space between paragraphs, no indent
\usepackage{microtype}      % subtle kerning/spacing improvements

% ---- Colors (muted palette, easy on light-sensitive eyes) -------------
\usepackage{xcolor}
\definecolor{pagebg}       {HTML}{FAF6F0}  % warm cream background
\definecolor{bodyfg}       {HTML}{3D3A36}  % warm dark gray body text
\definecolor{heading}      {HTML}{3D6A6B}  % muted teal
\definecolor{subheading}   {HTML}{5B7A9C}  % dusty blue
\definecolor{accent}       {HTML}{8A6B4A}  % warm tan
\definecolor{linkcol}      {HTML}{5B7A9C}  % same as subheading
\definecolor{mutedtext}    {HTML}{6B6560}  % secondary / caption text
\definecolor{rulecol}      {HTML}{C8BFB0}  % separator / rule

% Callout box palette
\definecolor{defnFrame}    {HTML}{9DB89B}  \definecolor{defnBack}    {HTML}{EEF2E8}
\definecolor{ideaFrame}    {HTML}{9FB4CD}  \definecolor{ideaBack}    {HTML}{E8EEF5}
\definecolor{warnFrame}    {HTML}{CFAE87}  \definecolor{warnBack}    {HTML}{F5EDDF}
\definecolor{noteFrame}    {HTML}{BFA6AE}  \definecolor{noteBack}    {HTML}{F2E8EE}
\definecolor{exFrame}      {HTML}{9FBDBD}  \definecolor{exBack}      {HTML}{E8F0F0}
\definecolor{codeBack}     {HTML}{F3EDE2}
\definecolor{codeBorder}   {HTML}{D4CBBA}
\definecolor{codeKeyword}  {HTML}{6B4F7A}
\definecolor{codeString}   {HTML}{8A6B4A}
\definecolor{codeComment}  {HTML}{8A857F}

% ---- Page background + base text color --------------------------------
\usepackage{pagecolor}
\pagecolor{pagebg}
\color{bodyfg}

% ---- Fonts ------------------------------------------------------------
\usepackage[T1]{fontenc}
\IfFileExists{libertine.sty}{%
    \usepackage{libertine}      % warm, readable serif
    \usepackage[scaled=0.95]{inconsolata}  % softer monospace
}{}

% ---- Hyperref ---------------------------------------------------------
\usepackage{hyperref}
\hypersetup{
    colorlinks=true,
    linkcolor=linkcol,
    urlcolor=linkcol,
    citecolor=linkcol,
    pdfborder={0 0 0},
}

% ---- Section styling --------------------------------------------------
\usepackage[explicit]{titlesec}
\titleformat{\section}
    {\color{heading}\Large\bfseries}
    {\color{heading}\thesection\hspace{0.6em}}{0pt}{#1}
\titleformat{\subsection}
    {\color{subheading}\large\bfseries}
    {\color{subheading}\thesubsection\hspace{0.5em}}{0pt}{#1}
\titleformat{\subsubsection}
    {\color{accent}\normalsize\bfseries}
    {\color{accent}\thesubsubsection\hspace{0.5em}}{0pt}{#1}
\titleformat{\paragraph}[runin]
    {\color{accent}\bfseries}{}{0pt}{#1.}

\titlespacing*{\section}{0pt}{1.6em}{0.6em}
\titlespacing*{\subsection}{0pt}{1.2em}{0.4em}

% ---- Enumerate / itemize spacing --------------------------------------
\usepackage{enumitem}
\setlist[itemize]{leftmargin=1.4em, itemsep=2pt, topsep=4pt}
\setlist[enumerate]{leftmargin=1.6em, itemsep=2pt, topsep=4pt}

% ---- Code listings ----------------------------------------------------
\usepackage{listings}
\lstdefinestyle{notescode}{
    backgroundcolor=\color{codeBack},
    basicstyle=\ttfamily\small\color{bodyfg},
    keywordstyle=\color{codeKeyword}\bfseries,
    stringstyle=\color{codeString},
    commentstyle=\color{codeComment}\itshape,
    frame=single,
    framesep=6pt,
    rulecolor=\color{codeBorder},
    showstringspaces=false,
    breaklines=true,
    xleftmargin=10pt,
    xrightmargin=10pt,
    aboveskip=8pt,
    belowskip=8pt,
}
\lstset{style=notescode}

% ---- Callout boxes (tcolorbox) ----------------------------------------
\usepackage[most]{tcolorbox}

\newtcolorbox{defnbox}[1][]{
    enhanced, breakable,
    colback=defnBack, colframe=defnFrame, coltitle=bodyfg,
    title={\bfseries Definition \textemdash\ #1},
    fonttitle=\bfseries,
    boxrule=0.5pt, arc=2pt, left=8pt, right=8pt, top=6pt, bottom=6pt,
    attach boxed title to top left={xshift=10pt, yshift=-6pt},
    boxed title style={colback=defnFrame, colframe=defnFrame, arc=2pt},
    before skip=8pt, after skip=8pt,
}
\newtcolorbox{ideabox}[1][]{
    enhanced, breakable,
    colback=ideaBack, colframe=ideaFrame, coltitle=bodyfg,
    title={\bfseries Key idea #1},
    boxrule=0.5pt, arc=2pt, left=8pt, right=8pt, top=6pt, bottom=6pt,
    attach boxed title to top left={xshift=10pt, yshift=-6pt},
    boxed title style={colback=ideaFrame, colframe=ideaFrame, arc=2pt},
    before skip=8pt, after skip=8pt,
}
\newtcolorbox{warnbox}[1][]{
    enhanced, breakable,
    colback=warnBack, colframe=warnFrame, coltitle=bodyfg,
    title={\bfseries Gotcha #1},
    boxrule=0.5pt, arc=2pt, left=8pt, right=8pt, top=6pt, bottom=6pt,
    attach boxed title to top left={xshift=10pt, yshift=-6pt},
    boxed title style={colback=warnFrame, colframe=warnFrame, arc=2pt},
    before skip=8pt, after skip=8pt,
}
\newtcolorbox{notebox}[1][]{
    enhanced, breakable,
    colback=noteBack, colframe=noteFrame, coltitle=bodyfg,
    title={\bfseries Aside #1},
    boxrule=0.5pt, arc=2pt, left=8pt, right=8pt, top=6pt, bottom=6pt,
    attach boxed title to top left={xshift=10pt, yshift=-6pt},
    boxed title style={colback=noteFrame, colframe=noteFrame, arc=2pt},
    before skip=8pt, after skip=8pt,
}
\newtcolorbox{examplebox}[1][]{
    enhanced, breakable,
    colback=exBack, colframe=exFrame, coltitle=bodyfg,
    title={\bfseries Example #1},
    boxrule=0.5pt, arc=2pt, left=8pt, right=8pt, top=6pt, bottom=6pt,
    attach boxed title to top left={xshift=10pt, yshift=-6pt},
    boxed title style={colback=exFrame, colframe=exFrame, arc=2pt},
    before skip=8pt, after skip=8pt,
}

% ---- TikZ graph styles (added 2026-04-25 for ch_10 augmentation) ------
% tcolorbox[most] already loads tikz; this block declares the libraries
% + shared `vertex` / `edge` styles used by ch_10's general directed-graph
% diagrams. ch_9's tree-style TikZ uses inline overrides and is
% unaffected. Simple circle nodes + arrow edges only -- no production
% styling.
\usetikzlibrary{arrows.meta, positioning, calc}
\tikzset{
    vertex/.style={circle, draw, minimum size=7mm, inner sep=0pt,
                   font=\small, thick},
    vsource/.style={vertex, fill=ideaBack},
    vdone/.style={vertex, fill=defnBack},
    edge/.style={->, >=Stealth, thick},
    uedge/.style={-, thick},
    treeedge/.style={->, >=Stealth, thick},
    backedge/.style={->, >=Stealth, thick, dashed, red!70!black},
    fwdedge/.style={->, >=Stealth, thick, blue!60!black},
    crossedge/.style={->, >=Stealth, thick, green!50!black, dotted},
    mstedge/.style={-, very thick, accent},
}

% ---- Title block ------------------------------------------------------
\usepackage{titling}
\pretitle{\begin{center}\color{heading}\LARGE\bfseries}
\posttitle{\end{center}\vskip -0.4em}
\preauthor{\begin{center}\color{mutedtext}\normalsize}
\postauthor{\end{center}}
\predate{\begin{center}\color{mutedtext}\small}
\postdate{\end{center}\vspace{0.4em}{\color{rulecol}\hrule height 0.4pt}\vspace{1.2em}}
` right after \documentclass
%   in each chapter's lectures.tex and notes.tex.
% =======================================================================

% ---- Layout -----------------------------------------------------------
\usepackage[margin=1in]{geometry}
\usepackage{parskip}        % space between paragraphs, no indent
\usepackage{microtype}      % subtle kerning/spacing improvements

% ---- Colors (muted palette, easy on light-sensitive eyes) -------------
\usepackage{xcolor}
\definecolor{pagebg}       {HTML}{FAF6F0}  % warm cream background
\definecolor{bodyfg}       {HTML}{3D3A36}  % warm dark gray body text
\definecolor{heading}      {HTML}{3D6A6B}  % muted teal
\definecolor{subheading}   {HTML}{5B7A9C}  % dusty blue
\definecolor{accent}       {HTML}{8A6B4A}  % warm tan
\definecolor{linkcol}      {HTML}{5B7A9C}  % same as subheading
\definecolor{mutedtext}    {HTML}{6B6560}  % secondary / caption text
\definecolor{rulecol}      {HTML}{C8BFB0}  % separator / rule

% Callout box palette
\definecolor{defnFrame}    {HTML}{9DB89B}  \definecolor{defnBack}    {HTML}{EEF2E8}
\definecolor{ideaFrame}    {HTML}{9FB4CD}  \definecolor{ideaBack}    {HTML}{E8EEF5}
\definecolor{warnFrame}    {HTML}{CFAE87}  \definecolor{warnBack}    {HTML}{F5EDDF}
\definecolor{noteFrame}    {HTML}{BFA6AE}  \definecolor{noteBack}    {HTML}{F2E8EE}
\definecolor{exFrame}      {HTML}{9FBDBD}  \definecolor{exBack}      {HTML}{E8F0F0}
\definecolor{codeBack}     {HTML}{F3EDE2}
\definecolor{codeBorder}   {HTML}{D4CBBA}
\definecolor{codeKeyword}  {HTML}{6B4F7A}
\definecolor{codeString}   {HTML}{8A6B4A}
\definecolor{codeComment}  {HTML}{8A857F}

% ---- Page background + base text color --------------------------------
\usepackage{pagecolor}
\pagecolor{pagebg}
\color{bodyfg}

% ---- Fonts ------------------------------------------------------------
\usepackage[T1]{fontenc}
\IfFileExists{libertine.sty}{%
    \usepackage{libertine}      % warm, readable serif
    \usepackage[scaled=0.95]{inconsolata}  % softer monospace
}{}

% ---- Hyperref ---------------------------------------------------------
\usepackage{hyperref}
\hypersetup{
    colorlinks=true,
    linkcolor=linkcol,
    urlcolor=linkcol,
    citecolor=linkcol,
    pdfborder={0 0 0},
}

% ---- Section styling --------------------------------------------------
\usepackage[explicit]{titlesec}
\titleformat{\section}
    {\color{heading}\Large\bfseries}
    {\color{heading}\thesection\hspace{0.6em}}{0pt}{#1}
\titleformat{\subsection}
    {\color{subheading}\large\bfseries}
    {\color{subheading}\thesubsection\hspace{0.5em}}{0pt}{#1}
\titleformat{\subsubsection}
    {\color{accent}\normalsize\bfseries}
    {\color{accent}\thesubsubsection\hspace{0.5em}}{0pt}{#1}
\titleformat{\paragraph}[runin]
    {\color{accent}\bfseries}{}{0pt}{#1.}

\titlespacing*{\section}{0pt}{1.6em}{0.6em}
\titlespacing*{\subsection}{0pt}{1.2em}{0.4em}

% ---- Enumerate / itemize spacing --------------------------------------
\usepackage{enumitem}
\setlist[itemize]{leftmargin=1.4em, itemsep=2pt, topsep=4pt}
\setlist[enumerate]{leftmargin=1.6em, itemsep=2pt, topsep=4pt}

% ---- Code listings ----------------------------------------------------
\usepackage{listings}
\lstdefinestyle{notescode}{
    backgroundcolor=\color{codeBack},
    basicstyle=\ttfamily\small\color{bodyfg},
    keywordstyle=\color{codeKeyword}\bfseries,
    stringstyle=\color{codeString},
    commentstyle=\color{codeComment}\itshape,
    frame=single,
    framesep=6pt,
    rulecolor=\color{codeBorder},
    showstringspaces=false,
    breaklines=true,
    xleftmargin=10pt,
    xrightmargin=10pt,
    aboveskip=8pt,
    belowskip=8pt,
}
\lstset{style=notescode}

% ---- Callout boxes (tcolorbox) ----------------------------------------
\usepackage[most]{tcolorbox}

\newtcolorbox{defnbox}[1][]{
    enhanced, breakable,
    colback=defnBack, colframe=defnFrame, coltitle=bodyfg,
    title={\bfseries Definition \textemdash\ #1},
    fonttitle=\bfseries,
    boxrule=0.5pt, arc=2pt, left=8pt, right=8pt, top=6pt, bottom=6pt,
    attach boxed title to top left={xshift=10pt, yshift=-6pt},
    boxed title style={colback=defnFrame, colframe=defnFrame, arc=2pt},
    before skip=8pt, after skip=8pt,
}
\newtcolorbox{ideabox}[1][]{
    enhanced, breakable,
    colback=ideaBack, colframe=ideaFrame, coltitle=bodyfg,
    title={\bfseries Key idea #1},
    boxrule=0.5pt, arc=2pt, left=8pt, right=8pt, top=6pt, bottom=6pt,
    attach boxed title to top left={xshift=10pt, yshift=-6pt},
    boxed title style={colback=ideaFrame, colframe=ideaFrame, arc=2pt},
    before skip=8pt, after skip=8pt,
}
\newtcolorbox{warnbox}[1][]{
    enhanced, breakable,
    colback=warnBack, colframe=warnFrame, coltitle=bodyfg,
    title={\bfseries Gotcha #1},
    boxrule=0.5pt, arc=2pt, left=8pt, right=8pt, top=6pt, bottom=6pt,
    attach boxed title to top left={xshift=10pt, yshift=-6pt},
    boxed title style={colback=warnFrame, colframe=warnFrame, arc=2pt},
    before skip=8pt, after skip=8pt,
}
\newtcolorbox{notebox}[1][]{
    enhanced, breakable,
    colback=noteBack, colframe=noteFrame, coltitle=bodyfg,
    title={\bfseries Aside #1},
    boxrule=0.5pt, arc=2pt, left=8pt, right=8pt, top=6pt, bottom=6pt,
    attach boxed title to top left={xshift=10pt, yshift=-6pt},
    boxed title style={colback=noteFrame, colframe=noteFrame, arc=2pt},
    before skip=8pt, after skip=8pt,
}
\newtcolorbox{examplebox}[1][]{
    enhanced, breakable,
    colback=exBack, colframe=exFrame, coltitle=bodyfg,
    title={\bfseries Example #1},
    boxrule=0.5pt, arc=2pt, left=8pt, right=8pt, top=6pt, bottom=6pt,
    attach boxed title to top left={xshift=10pt, yshift=-6pt},
    boxed title style={colback=exFrame, colframe=exFrame, arc=2pt},
    before skip=8pt, after skip=8pt,
}

% ---- TikZ graph styles (added 2026-04-25 for ch_10 augmentation) ------
% tcolorbox[most] already loads tikz; this block declares the libraries
% + shared `vertex` / `edge` styles used by ch_10's general directed-graph
% diagrams. ch_9's tree-style TikZ uses inline overrides and is
% unaffected. Simple circle nodes + arrow edges only -- no production
% styling.
\usetikzlibrary{arrows.meta, positioning, calc}
\tikzset{
    vertex/.style={circle, draw, minimum size=7mm, inner sep=0pt,
                   font=\small, thick},
    vsource/.style={vertex, fill=ideaBack},
    vdone/.style={vertex, fill=defnBack},
    edge/.style={->, >=Stealth, thick},
    uedge/.style={-, thick},
    treeedge/.style={->, >=Stealth, thick},
    backedge/.style={->, >=Stealth, thick, dashed, red!70!black},
    fwdedge/.style={->, >=Stealth, thick, blue!60!black},
    crossedge/.style={->, >=Stealth, thick, green!50!black, dotted},
    mstedge/.style={-, very thick, accent},
}

% ---- Title block ------------------------------------------------------
\usepackage{titling}
\pretitle{\begin{center}\color{heading}\LARGE\bfseries}
\posttitle{\end{center}\vskip -0.4em}
\preauthor{\begin{center}\color{mutedtext}\normalsize}
\postauthor{\end{center}}
\predate{\begin{center}\color{mutedtext}\small}
\postdate{\end{center}\vspace{0.4em}{\color{rulecol}\hrule height 0.4pt}\vspace{1.2em}}
` right after \documentclass
%   in each chapter's lectures.tex and notes.tex.
% =======================================================================

% ---- Layout -----------------------------------------------------------
\usepackage[margin=1in]{geometry}
\usepackage{parskip}        % space between paragraphs, no indent
\usepackage{microtype}      % subtle kerning/spacing improvements

% ---- Colors (muted palette, easy on light-sensitive eyes) -------------
\usepackage{xcolor}
\definecolor{pagebg}       {HTML}{FAF6F0}  % warm cream background
\definecolor{bodyfg}       {HTML}{3D3A36}  % warm dark gray body text
\definecolor{heading}      {HTML}{3D6A6B}  % muted teal
\definecolor{subheading}   {HTML}{5B7A9C}  % dusty blue
\definecolor{accent}       {HTML}{8A6B4A}  % warm tan
\definecolor{linkcol}      {HTML}{5B7A9C}  % same as subheading
\definecolor{mutedtext}    {HTML}{6B6560}  % secondary / caption text
\definecolor{rulecol}      {HTML}{C8BFB0}  % separator / rule

% Callout box palette
\definecolor{defnFrame}    {HTML}{9DB89B}  \definecolor{defnBack}    {HTML}{EEF2E8}
\definecolor{ideaFrame}    {HTML}{9FB4CD}  \definecolor{ideaBack}    {HTML}{E8EEF5}
\definecolor{warnFrame}    {HTML}{CFAE87}  \definecolor{warnBack}    {HTML}{F5EDDF}
\definecolor{noteFrame}    {HTML}{BFA6AE}  \definecolor{noteBack}    {HTML}{F2E8EE}
\definecolor{exFrame}      {HTML}{9FBDBD}  \definecolor{exBack}      {HTML}{E8F0F0}
\definecolor{codeBack}     {HTML}{F3EDE2}
\definecolor{codeBorder}   {HTML}{D4CBBA}
\definecolor{codeKeyword}  {HTML}{6B4F7A}
\definecolor{codeString}   {HTML}{8A6B4A}
\definecolor{codeComment}  {HTML}{8A857F}

% ---- Page background + base text color --------------------------------
\usepackage{pagecolor}
\pagecolor{pagebg}
\color{bodyfg}

% ---- Fonts ------------------------------------------------------------
\usepackage[T1]{fontenc}
\IfFileExists{libertine.sty}{%
    \usepackage{libertine}      % warm, readable serif
    \usepackage[scaled=0.95]{inconsolata}  % softer monospace
}{}

% ---- Hyperref ---------------------------------------------------------
\usepackage{hyperref}
\hypersetup{
    colorlinks=true,
    linkcolor=linkcol,
    urlcolor=linkcol,
    citecolor=linkcol,
    pdfborder={0 0 0},
}

% ---- Section styling --------------------------------------------------
\usepackage[explicit]{titlesec}
\titleformat{\section}
    {\color{heading}\Large\bfseries}
    {\color{heading}\thesection\hspace{0.6em}}{0pt}{#1}
\titleformat{\subsection}
    {\color{subheading}\large\bfseries}
    {\color{subheading}\thesubsection\hspace{0.5em}}{0pt}{#1}
\titleformat{\subsubsection}
    {\color{accent}\normalsize\bfseries}
    {\color{accent}\thesubsubsection\hspace{0.5em}}{0pt}{#1}
\titleformat{\paragraph}[runin]
    {\color{accent}\bfseries}{}{0pt}{#1.}

\titlespacing*{\section}{0pt}{1.6em}{0.6em}
\titlespacing*{\subsection}{0pt}{1.2em}{0.4em}

% ---- Enumerate / itemize spacing --------------------------------------
\usepackage{enumitem}
\setlist[itemize]{leftmargin=1.4em, itemsep=2pt, topsep=4pt}
\setlist[enumerate]{leftmargin=1.6em, itemsep=2pt, topsep=4pt}

% ---- Code listings ----------------------------------------------------
\usepackage{listings}
\lstdefinestyle{notescode}{
    backgroundcolor=\color{codeBack},
    basicstyle=\ttfamily\small\color{bodyfg},
    keywordstyle=\color{codeKeyword}\bfseries,
    stringstyle=\color{codeString},
    commentstyle=\color{codeComment}\itshape,
    frame=single,
    framesep=6pt,
    rulecolor=\color{codeBorder},
    showstringspaces=false,
    breaklines=true,
    xleftmargin=10pt,
    xrightmargin=10pt,
    aboveskip=8pt,
    belowskip=8pt,
}
\lstset{style=notescode}

% ---- Callout boxes (tcolorbox) ----------------------------------------
\usepackage[most]{tcolorbox}

\newtcolorbox{defnbox}[1][]{
    enhanced, breakable,
    colback=defnBack, colframe=defnFrame, coltitle=bodyfg,
    title={\bfseries Definition \textemdash\ #1},
    fonttitle=\bfseries,
    boxrule=0.5pt, arc=2pt, left=8pt, right=8pt, top=6pt, bottom=6pt,
    attach boxed title to top left={xshift=10pt, yshift=-6pt},
    boxed title style={colback=defnFrame, colframe=defnFrame, arc=2pt},
    before skip=8pt, after skip=8pt,
}
\newtcolorbox{ideabox}[1][]{
    enhanced, breakable,
    colback=ideaBack, colframe=ideaFrame, coltitle=bodyfg,
    title={\bfseries Key idea #1},
    boxrule=0.5pt, arc=2pt, left=8pt, right=8pt, top=6pt, bottom=6pt,
    attach boxed title to top left={xshift=10pt, yshift=-6pt},
    boxed title style={colback=ideaFrame, colframe=ideaFrame, arc=2pt},
    before skip=8pt, after skip=8pt,
}
\newtcolorbox{warnbox}[1][]{
    enhanced, breakable,
    colback=warnBack, colframe=warnFrame, coltitle=bodyfg,
    title={\bfseries Gotcha #1},
    boxrule=0.5pt, arc=2pt, left=8pt, right=8pt, top=6pt, bottom=6pt,
    attach boxed title to top left={xshift=10pt, yshift=-6pt},
    boxed title style={colback=warnFrame, colframe=warnFrame, arc=2pt},
    before skip=8pt, after skip=8pt,
}
\newtcolorbox{notebox}[1][]{
    enhanced, breakable,
    colback=noteBack, colframe=noteFrame, coltitle=bodyfg,
    title={\bfseries Aside #1},
    boxrule=0.5pt, arc=2pt, left=8pt, right=8pt, top=6pt, bottom=6pt,
    attach boxed title to top left={xshift=10pt, yshift=-6pt},
    boxed title style={colback=noteFrame, colframe=noteFrame, arc=2pt},
    before skip=8pt, after skip=8pt,
}
\newtcolorbox{examplebox}[1][]{
    enhanced, breakable,
    colback=exBack, colframe=exFrame, coltitle=bodyfg,
    title={\bfseries Example #1},
    boxrule=0.5pt, arc=2pt, left=8pt, right=8pt, top=6pt, bottom=6pt,
    attach boxed title to top left={xshift=10pt, yshift=-6pt},
    boxed title style={colback=exFrame, colframe=exFrame, arc=2pt},
    before skip=8pt, after skip=8pt,
}

% ---- TikZ graph styles (added 2026-04-25 for ch_10 augmentation) ------
% tcolorbox[most] already loads tikz; this block declares the libraries
% + shared `vertex` / `edge` styles used by ch_10's general directed-graph
% diagrams. ch_9's tree-style TikZ uses inline overrides and is
% unaffected. Simple circle nodes + arrow edges only -- no production
% styling.
\usetikzlibrary{arrows.meta, positioning, calc}
\tikzset{
    vertex/.style={circle, draw, minimum size=7mm, inner sep=0pt,
                   font=\small, thick},
    vsource/.style={vertex, fill=ideaBack},
    vdone/.style={vertex, fill=defnBack},
    edge/.style={->, >=Stealth, thick},
    uedge/.style={-, thick},
    treeedge/.style={->, >=Stealth, thick},
    backedge/.style={->, >=Stealth, thick, dashed, red!70!black},
    fwdedge/.style={->, >=Stealth, thick, blue!60!black},
    crossedge/.style={->, >=Stealth, thick, green!50!black, dotted},
    mstedge/.style={-, very thick, accent},
}

% ---- Title block ------------------------------------------------------
\usepackage{titling}
\pretitle{\begin{center}\color{heading}\LARGE\bfseries}
\posttitle{\end{center}\vskip -0.4em}
\preauthor{\begin{center}\color{mutedtext}\normalsize}
\postauthor{\end{center}}
\predate{\begin{center}\color{mutedtext}\small}
\postdate{\end{center}\vspace{0.4em}{\color{rulecol}\hrule height 0.4pt}\vspace{1.2em}}


\title{CS 300 -- Chapter 4 Lectures\\\large Lists, Stacks, Queues, and Deques}
\author{Jose de Lima}
\date{\today}

\begin{document}
\maketitle

\begin{ideabox}[Chapter map]
\textbf{Where this sits in CS 300.} Chapter 3 named the seven data
structures; chapter 4 picks one family -- \emph{linear, ordered
collections} -- and builds it end-to-end, twice (pointer-based and
array-based). Everything you learn here about \emph{cost is a function
of memory layout} transfers verbatim to trees (ch.~6), heaps, and graph
adjacency lists (opt.\ ch.~10).

\textbf{What you add to your toolkit here:}
\begin{itemize}
    \item \textbf{Pointer-based lists.} Singly linked (SLL) and doubly
          linked (DLL); append, prepend, insert-after, remove, search,
          traverse. You can write these from scratch, by hand.
    \item \textbf{Sentinel / dummy-node idiom.} The trick that collapses
          ``first / middle / last'' special cases into one uniform case.
    \item \textbf{Three restricted list ADTs}: Stack (LIFO), Queue
          (FIFO), Deque (both ends) -- each built on linked lists or
          arrays.
    \item \textbf{Array-based lists.} The \texttt{std::vector} cost
          model (amortized $O(1)$ back, $O(n)$ insert/erase elsewhere)
          contrasted cleanly with the linked-list cost model.
    \item \textbf{A decision framework} (see 4.18+) for picking between
          array-based, SLL, DLL, and sentinel-DLL. This is the question
          that actually shows up on assignments.
\end{itemize}

\textbf{Mastery checklist} -- you should be able to, without references:
\begin{enumerate}
    \item Draw an SLL of 3 nodes and execute \texttt{insertAfter} and
          \texttt{remove} on it by redrawing the pointers one step at a
          time.
    \item Write the node \texttt{struct} for an SLL and a DLL; write
          \texttt{append}, \texttt{prepend}, \texttt{insertAfter},
          \texttt{remove}, and \texttt{search} for both.
    \item Explain why \texttt{remove} on an SLL requires the
          \emph{predecessor}'s pointer, and how DLL gets around it.
    \item State the cost of every List-ADT operation on (a) SLL,
          (b) DLL, (c) array-based list -- without guessing. Know
          \emph{why} each one is what it is.
    \item Implement Stack, Queue, and Deque on both a linked list and
          an array, and argue which backing data structure you'd pick
          for each.
    \item Rewrite an SLL method that has a 3-case switch (empty / head /
          middle) to use a sentinel node and collapse to one case.
    \item Traverse a linked list iteratively \emph{and} recursively;
          know when recursion blows the stack.
    \item Given a workload description (``mostly appends, rare middle
          removals, needs $O(1)$ random access''), pick the right list
          variant and justify.
\end{enumerate}

\textbf{Looking ahead.} Chapter 5 moves from linear to keyed storage
(hash tables), but the pointer-chasing mental model you're building here
is exactly what hash-table chaining uses internally. Chapter 6 (trees)
extends the node + pointer idiom from one ``next'' pointer to two
(``left'' / ``right''); the whole tree chapter will feel like a
generalization of this one.
\end{ideabox}

\section{4.1 The List ADT}

Chapter 3 surveyed the seven fundamental data structures and the
nine common abstract data types, then spent most of its energy on
searching and sorting. Chapter 4 starts the depth-first tour: each
ADT gets its own chapter where we define the contract, implement it,
and analyze the operations. First up: the \textbf{list}.

The list is the workhorse ADT. Nearly every other sequence-based
structure (stack, queue, deque) is a list with additional
constraints. Understanding what a list's contract actually
\emph{is} -- and what it deliberately leaves unspecified -- is the
first step to picking the right structure later.

\begin{defnbox}[List]
A \textbf{list} is an abstract data type for a \emph{finite, ordered
sequence} of items. Items have \emph{positions} (first, second, last),
duplicates are allowed, and operations are defined in terms of those
positions: append at the end, prepend at the front, insert after a
given item, remove, search, iterate.
\end{defnbox}

\begin{ideabox}[Three things the list ADT guarantees]
\begin{enumerate}[leftmargin=*]
  \item \textbf{Order matters.} \texttt{(7, 9, 5)} is different from
        \texttt{(5, 7, 9)}. A list is not a set.
  \item \textbf{Duplicates are allowed.} \texttt{(3, 3, 3)} is a
        valid three-item list.
  \item \textbf{Positions are stable.} ``The second item'' refers to
        a specific slot, not just ``any item.'' Inserting at the
        front shifts every other item's position index by one;
        the ADT is aware of this.
\end{enumerate}
\end{ideabox}

\subsection{The contract: operations}

\begin{center}
\footnotesize
\renewcommand{\arraystretch}{1.3}
\begin{tabular}{@{}p{3.4cm}p{5cm}p{4.8cm}@{}}
\hline
\textbf{Operation} & \textbf{What it does} & \textbf{Example (start: 99, 77)} \\
\hline
\texttt{Append(L, x)}          & Insert $x$ at the end              & \texttt{Append(L, 44)} $\Rightarrow$ 99, 77, 44 \\
\texttt{Prepend(L, x)}         & Insert $x$ at the front            & \texttt{Prepend(L, 44)} $\Rightarrow$ 44, 99, 77 \\
\texttt{InsertAfter(L, w, x)}  & Insert $x$ immediately after $w$   & \texttt{InsertAfter(L, 99, 44)} $\Rightarrow$ 99, 44, 77 \\
\texttt{Remove(L, x)}          & Remove first occurrence of $x$     & \texttt{Remove(L, 77)} $\Rightarrow$ 99 \\
\texttt{Search(L, x)}          & Return item (or \texttt{null})     & \texttt{Search(L, 99)} $\Rightarrow$ 99; \texttt{Search(L, 22)} $\Rightarrow$ null \\
\texttt{Print(L)}              & Print items in order               & ``99, 77'' \\
\texttt{PrintReverse(L)}       & Print items in reverse order       & ``77, 99'' \\
\texttt{Sort(L)}               & Sort items ascending               & 77, 99 \\
\texttt{IsEmpty(L)}            & True iff list has no items         & \texttt{false} for (99, 77) \\
\texttt{GetLength(L)}          & Number of items                    & 2 \\
\hline
\end{tabular}
\end{center}

\begin{notebox}[Not every list exposes every operation]
The table above is a \emph{menu}, not a mandatory interface. Real
list libraries pick a subset. \texttt{std::vector} exposes indexed
access (\texttt{v[i]}) but no constant-time \texttt{Prepend};
\texttt{std::forward\_list} exposes \texttt{Prepend} (via
\texttt{push\_front}) but not indexed access. The ADT names the
operations; the implementation picks which to make cheap.
\end{notebox}

\subsection{What the contract \emph{doesn't} say}

\begin{ideabox}[Conspicuous silences]
The list ADT says nothing about:
\begin{itemize}[leftmargin=*]
  \item \textbf{How items are stored in memory.} Contiguously?
        Linked nodes? Chunks? All valid.
  \item \textbf{Cost of operations.} The ADT defines what \texttt{Append}
        \emph{means}, not how fast it runs. An array-backed list and
        a singly-linked list both implement \texttt{Append}; they
        pay very different prices for it.
  \item \textbf{Indexed random access.} The raw list ADT is defined
        by position-relative operations (``after this item,''
        ``first,'' ``last''), not by integer index. The \emph{dynamic
        array} ADT is the refinement that adds $O(1)$
        \texttt{at(i)}.
\end{itemize}
The silences are the point. Leaving implementation unspecified means
callers can substitute one list implementation for another without
rewriting application code.
\end{ideabox}

\begin{notebox}[Two interface families: Sequence and Set]
Step back one level. Every container ADT in CS 300 belongs to one of
two families, distinguished by where the order of items lives:
\begin{itemize}[leftmargin=*]
  \item \textbf{Sequence interface.} Items have an \emph{extrinsic}
        order --- the caller chose where each item goes. Operations
        are positional: \texttt{get\_at(i)}, \texttt{set\_at(i,x)},
        \texttt{insert\_at(i,x)}, \texttt{insert\_first/last(x)},
        \texttt{delete\_first/last()}. \emph{This chapter implements
        Sequence.} List, stack, queue, and deque are all Sequence
        with restricted operation sets.
  \item \textbf{Set interface.} Items have an \emph{intrinsic} order
        determined by a key. Operations are key-based:
        \texttt{find(k)}, \texttt{insert(x)}, \texttt{delete(k)},
        and ordered ops like \texttt{find\_min()},
        \texttt{find\_next(k)}. ch.~5 (hash tables) and ch.~6 (BSTs)
        implement Set.
\end{itemize}
The framing is from MIT 6.006. It explains why the chapter map is the
shape it is: ch.~4 builds Sequence implementations end-to-end, then
ch.~5--ch.~6 swap the interface and rebuild from scratch. Pick the
right family first; pick the implementation second.
\end{notebox}

\subsection{List is the root ADT; stacks and queues restrict it}

\begin{defnbox}[Restricted-access lists]
Several ADTs later in the chapter are just lists with
\emph{restricted} operations:
\begin{itemize}[leftmargin=*]
  \item \textbf{Stack}: a list where you can only insert or remove
        at \emph{one end} (the ``top''). LIFO.
  \item \textbf{Queue}: a list where you insert at one end and
        remove at the other. FIFO.
  \item \textbf{Deque}: insert or remove at \emph{either} end, but
        nowhere else.
\end{itemize}
Narrower interfaces, stronger guarantees about usage patterns,
faster implementations possible.
\end{defnbox}

\subsection{C++: \texttt{std::list}, \texttt{std::vector}, \texttt{std::forward\_list}}

C++ gives you several list-like containers. All three satisfy the
list ADT; they differ in which operations are cheap and what extra
features they offer.

\begin{examplebox}[Three takes on ``list'' in C++]
\begin{lstlisting}[basicstyle=\ttfamily\small,frame=single,language=C++]
#include <vector>
#include <list>
#include <forward_list>

std::vector<int>        v = {7, 9, 5};   // dynamic array
std::list<int>          l = {7, 9, 5};   // doubly linked list
std::forward_list<int>  f = {7, 9, 5};   // singly linked list

v.push_back(44);                 // O(1) amortized (append)
l.push_back(44);                 // O(1)
// f has no push_back!           // forward_list is single-link only

v.insert(v.begin(), 44);         // O(n) (prepend: shifts everything)
l.push_front(44);                // O(1)
f.push_front(44);                // O(1)

int len_v = v.size();            // O(1)
int len_l = l.size();            // O(1)
// f.size() doesn't exist; forward_list doesn't track length
\end{lstlisting}
\end{examplebox}

\begin{center}
\footnotesize
\renewcommand{\arraystretch}{1.3}
\begin{tabular}{@{}lccc@{}}
\hline
\textbf{Operation}            & \texttt{std::vector} & \texttt{std::list} & \texttt{std::forward\_list} \\
\hline
Append at end                 & $O(1)$ amortized     & $O(1)$             & no direct op \\
Prepend at front              & $O(n)$               & $O(1)$             & $O(1)$ \\
Insert after a known position & $O(n)$               & $O(1)$             & $O(1)$ \\
Indexed access \texttt{v[i]}  & $O(1)$               & $O(i)$             & $O(i)$ \\
Search by value               & $O(n)$               & $O(n)$             & $O(n)$ \\
\texttt{size()}               & $O(1)$               & $O(1)$             & not provided \\
\hline
\end{tabular}
\end{center}

\begin{warnbox}[Pick the implementation from the \emph{access pattern}]
``I need a list'' doesn't tell you which container to use. Ask:
\begin{itemize}[leftmargin=*]
  \item Will I mostly \texttt{push\_back}? \texttt{vector}.
  \item Will I insert and remove in the middle a lot? \texttt{list}.
  \item Do I need indexed random access? \texttt{vector}.
  \item Do I need to splice or move sublists around? \texttt{list}.
\end{itemize}
For 80\% of everyday code, \texttt{std::vector} wins even when it
looks like the ``wrong'' choice, because its cache behavior and
constant-factor speed beat the linked-list options for moderate $n$.
\end{warnbox}

\subsection{What this chapter builds}

\begin{ideabox}[The plan]
\begin{itemize}[leftmargin=*]
  \item \textbf{Singly linked list.} Nodes with one \texttt{next}
        pointer. Implements the list ADT with $O(1)$ insert at
        known positions but no backward traversal.
  \item \textbf{Doubly linked list.} Nodes with \texttt{next} and
        \texttt{prev} pointers. Enables $O(1)$ removal given a
        node pointer, and efficient reverse iteration.
  \item \textbf{Array lists (dynamic arrays).} Contiguous storage,
        $O(1)$ indexed access, amortized $O(1)$ \texttt{push\_back}.
  \item \textbf{Stacks and queues.} Specializations of the above.
\end{itemize}
By the end of the chapter, every row in the C++ table above will be
something you can explain, not just memorize.
\end{ideabox}

\begin{notebox}[The pattern we'll repeat from now on]
For every ADT in this chapter and the next, we'll do:
\begin{enumerate}[leftmargin=*]
  \item Define the ADT (today).
  \item Implement it with one or two structures.
  \item Analyze each operation's cost.
  \item Compare to the alternatives.
\end{enumerate}
You saw this cycle informally for the sorts in Chapter 3. From here
on, it's the standard structure of every chapter.
\end{notebox}

\section{4.2 Singly Linked Lists: Append and Prepend}

This is the first concrete implementation of the list ADT: the
\textbf{singly linked list}. Instead of storing items in a contiguous
block of memory (like an array), a linked list scatters them across
individually-allocated \textbf{nodes} that ``point'' to each other.
The payoff is cheap insertion and deletion at known positions; the
price is slower traversal and higher per-item memory overhead.

\begin{defnbox}[Singly linked list]
A \textbf{singly linked list} is a sequence of nodes, where each
\textbf{node} stores a data value and a \texttt{next} pointer to the
following node. The list itself holds a \textbf{head} pointer (to
the first node) and usually a \textbf{tail} pointer (to the last
node). The last node's \texttt{next} is \texttt{nullptr}.
\end{defnbox}

\begin{ideabox}[Picture the list]
\texttt{head $\to$ [7 | $\bullet$] $\to$ [9 | $\bullet$] $\to$ [5 | null] $\leftarrow$ tail}

Each bracketed cell is a node. The left half is the data; the right
half is the \texttt{next} pointer. \texttt{head} points to the first
cell; \texttt{tail} points to the last; the last cell's \texttt{next}
is \texttt{nullptr}.
\end{ideabox}

\subsection{The node type in C++}

\begin{examplebox}[A node and a list wrapper]
\begin{lstlisting}[basicstyle=\ttfamily\small,frame=single,language=C++]
struct Node {
    int   data;
    Node* next;

    Node(int d) : data(d), next(nullptr) {}
};

struct LinkedList {
    Node* head = nullptr;
    Node* tail = nullptr;
};
\end{lstlisting}
\texttt{head} is how the list knows where to start; \texttt{tail} is
the optimization that lets \texttt{append} run in $O(1)$ without
walking the list. An empty list has both set to \texttt{nullptr}.
\end{examplebox}

\begin{warnbox}[Why \texttt{tail} is load-bearing]
Without a \texttt{tail} pointer, appending to a linked list would
require walking from \texttt{head} to the end -- $O(n)$ per append,
$O(n^2)$ to build the list. \emph{Always} cache the tail when you
care about append speed.
\end{warnbox}

\subsection{Append: add at the end}

Two cases: an empty list needs its \texttt{head} and \texttt{tail}
both pointed at the new node; a non-empty list only has to reach
through the existing tail.

\begin{examplebox}[\texttt{append}]
\begin{lstlisting}[basicstyle=\ttfamily\small,frame=single,language=C++]
void append(LinkedList& list, Node* newNode) {
    if (list.head == nullptr) {          // empty list
        list.head = newNode;
        list.tail = newNode;
    } else {                              // non-empty list
        list.tail->next = newNode;        // stitch on the end
        list.tail       = newNode;        // advance tail
    }
}
\end{lstlisting}
\end{examplebox}

\begin{ideabox}[Why \texttt{append} is $O(1)$]
The work is constant: check one pointer, write two pointers. No
scan of the list. The pre-existing \texttt{newNode$\to$next} is
assumed to already be \texttt{nullptr} from its constructor; if
you're reusing a node, set it explicitly.
\end{ideabox}

\begin{warnbox}[The empty-list branch is not optional]
If you forget the empty-list check and write
\texttt{list.tail$\to$next = newNode} when \texttt{tail} is
\texttt{nullptr}, you get a null-pointer dereference and a crash.
Every linked-list operation that touches \texttt{head} or
\texttt{tail} needs to handle the empty case explicitly. This is
the single most common linked-list bug.
\end{warnbox}

\subsection{Prepend: add at the front}

Symmetric to append, but we touch \texttt{head} instead of
\texttt{tail}.

\begin{examplebox}[\texttt{prepend}]
\begin{lstlisting}[basicstyle=\ttfamily\small,frame=single,language=C++]
void prepend(LinkedList& list, Node* newNode) {
    if (list.head == nullptr) {          // empty list
        list.head = newNode;
        list.tail = newNode;
    } else {                              // non-empty list
        newNode->next = list.head;        // point new node at old head
        list.head     = newNode;          // new node is the new head
    }
}
\end{lstlisting}
\end{examplebox}

\begin{ideabox}[The order of the two assignments matters]
In the non-empty branch, \texttt{newNode$\to$next = list.head} must
come \emph{before} \texttt{list.head = newNode}. Swap them and you
overwrite \texttt{head} with \texttt{newNode} first -- then the
subsequent \texttt{newNode$\to$next = list.head} makes the node
point at itself, losing the rest of the list. This is the classic
pointer-update gotcha: \textbf{always copy out before you overwrite}.
\end{ideabox}

\begin{notebox}[Why prepending into an array is different]
On an array-backed list, prepending means shifting every existing
item one slot to the right: $O(n)$. On a linked list it's just two
pointer writes: $O(1)$. The linked list doesn't have to keep items
adjacent in memory, so it never needs to shift.
\end{notebox}

\subsection{Cost summary}

\begin{center}
\footnotesize
\renewcommand{\arraystretch}{1.3}
\begin{tabular}{@{}lll@{}}
\hline
\textbf{Operation}    & \textbf{Singly linked list} & \textbf{Why} \\
\hline
\texttt{append}       & $O(1)$ with tail pointer    & Two writes; no traversal \\
\texttt{prepend}      & $O(1)$                      & Two writes; no traversal \\
\hline
\end{tabular}
\end{center}

Both basic insertions are constant time. That's the primary
advantage of a linked list over an array for workloads that
repeatedly add at the ends.

\subsection{Language conventions for ``null''}

\begin{defnbox}[Null]
A \textbf{null pointer} is a pointer that points to nothing. Every
pointer must either point to a valid object or be explicitly null.
Languages name this differently:
\begin{itemize}[leftmargin=*]
  \item C++: \texttt{nullptr} (modern), \texttt{NULL} (old C-style macro, avoid)
  \item C: \texttt{NULL}
  \item Java: \texttt{null}
  \item Python: \texttt{None}
  \item Rust / Swift: they avoid null via the \texttt{Option} type
\end{itemize}
In all cases the concept is the same: a sentinel value indicating
``no object here.''
\end{defnbox}

\begin{warnbox}[Use \texttt{nullptr}, not \texttt{NULL} or \texttt{0}]
In modern C++ (C++11 and later), always use \texttt{nullptr}. It's
type-safe: \texttt{NULL} is often defined as plain \texttt{0}, which
can accidentally be taken as an integer in overload resolution.
\texttt{nullptr} has its own dedicated type (\texttt{nullptr\_t}) and
can't be confused with an integer. Every linked-list node in this
course will initialize its pointers to \texttt{nullptr}.
\end{warnbox}

\subsection{What's coming next}

\begin{ideabox}[The remaining SLL operations]
Append and prepend are the easy cases because they touch the list's
endpoints, where we already have pointers. The operations coming in
4.3 -- insert-after-a-known-node, and remove-by-value -- are harder
because they involve traversing the list or keeping a \emph{trailing}
pointer to the node \emph{before} the one you want to modify. That
trailing-pointer pattern is the core awkwardness of singly linked
lists and the main motivation for \textbf{doubly} linked lists a
few sections later.
\end{ideabox}

\begin{notebox}[C++ owns its memory explicitly]
The pseudocode takes a \texttt{newNode} as a parameter and doesn't
discuss where it came from. In C++, you'll usually \texttt{new Node(x)}
to allocate and \texttt{delete} to free. Later, \texttt{std::list}
does all of that inside its \texttt{push\_back(x)}. This chapter
works with raw pointers to expose the mechanics; in production code
you'd use the STL container or smart pointers (\texttt{std::unique\_ptr})
to automate the memory management.
\end{notebox}

\section{4.3 Singly Linked Lists: Insert-After}

Append and prepend covered the two endpoints. \textbf{Insert-after}
is the general case: given a pointer to some existing node
\texttt{curNode}, splice a \texttt{newNode} in right after it. The
operation is still $O(1)$ -- because we already have a pointer to
the splice point -- but the bookkeeping has to handle three cases
instead of two.

\begin{defnbox}[Insert-after]
\textbf{InsertAfter(list, curNode, newNode)} inserts \texttt{newNode}
immediately after \texttt{curNode} in the list. If \texttt{curNode}
is \texttt{null} (or the list is empty), \texttt{newNode} becomes
the only node. If \texttt{curNode} is the current \texttt{tail},
\texttt{newNode} becomes the new tail.
\end{defnbox}

\subsection{The three cases}

The pseudocode looks like it's handling three unrelated scenarios,
but they all flow from the same question: \emph{which of the list's
invariants does this insertion threaten?}

\begin{ideabox}[What the list has to stay true to]
A well-formed singly linked list always satisfies:
\begin{enumerate}[leftmargin=*]
  \item \texttt{head} points to the first node or is \texttt{null}.
  \item \texttt{tail} points to the last node or is \texttt{null}.
  \item The last node's \texttt{next} is \texttt{null}.
  \item \texttt{head} is \texttt{null} iff \texttt{tail} is \texttt{null}.
\end{enumerate}
Every insertion must preserve all four. The three cases in the
algorithm each exist because a different invariant is at risk.
\end{ideabox}

\begin{center}
\footnotesize
\renewcommand{\arraystretch}{1.3}
\begin{tabular}{@{}p{3.8cm}p{8.2cm}@{}}
\hline
\textbf{Case} & \textbf{What needs updating} \\
\hline
Empty list                    & Both \texttt{head} and \texttt{tail} \\
\texttt{curNode} is the tail  & Tail's \texttt{next} pointer, and \texttt{tail} itself \\
\texttt{curNode} is interior  & Just two \texttt{next} pointers (no head/tail changes) \\
\hline
\end{tabular}
\end{center}

\subsection{The algorithm}

\begin{examplebox}[\texttt{insertAfter}]
\begin{lstlisting}[basicstyle=\ttfamily\small,frame=single,language=C++]
void insertAfter(LinkedList& list, Node* curNode, Node* newNode) {
    if (list.head == nullptr) {                  // Case 1: empty list
        list.head = newNode;
        list.tail = newNode;
    }
    else if (curNode == list.tail) {              // Case 2: after tail
        list.tail->next = newNode;
        list.tail       = newNode;
    }
    else {                                         // Case 3: interior
        newNode->next = curNode->next;
        curNode->next = newNode;
    }
}
\end{lstlisting}
\end{examplebox}

\begin{warnbox}[Case 3 is where most beginners trip]
The interior case has two assignments, and \emph{order matters}:
\begin{enumerate}[leftmargin=*]
  \item \texttt{newNode$\to$next = curNode$\to$next} -- make
        \texttt{newNode} point at what used to come after
        \texttt{curNode}. \emph{Copy out before overwriting.}
  \item \texttt{curNode$\to$next = newNode} -- then slot
        \texttt{newNode} in after \texttt{curNode}.
\end{enumerate}
Swap them and you overwrite \texttt{curNode$\to$next} first, losing
the rest of the list forever. Same pointer-update discipline as in
prepend -- make sure the new node is fully stitched \emph{before}
you change the node already in the list.
\end{warnbox}

\subsection*{Picture of the interior splice}

\begin{center}
\footnotesize
Before:\\
\texttt{... $\to$ [curNode | $\bullet$] $\to$ [nextNode | ...] $\to$ ...}
\\[0.5em]
Step 1: \texttt{newNode$\to$next = curNode$\to$next}\\
\texttt{... $\to$ [curNode | $\bullet$] $\to$ [nextNode | ...] $\to$ ...}\\
\texttt{\phantom{... $\to$ [curNode | $\bullet$]\ }\ $\nwarrow$}\\
\texttt{\phantom{... $\to$ }[newNode | $\bullet$]}
\\[0.5em]
Step 2: \texttt{curNode$\to$next = newNode}\\
\texttt{... $\to$ [curNode | $\bullet$] $\to$ [newNode | $\bullet$] $\to$ [nextNode | ...] $\to$ ...}
\end{center}

Two pointer writes, no scan. Constant time.

\subsection{Why we need the curNode argument}

\begin{ideabox}[``Insert at position $k$'' isn't a list-ADT primitive]
A client that wants to insert at a specific numeric position $k$
has to \emph{walk} the list to find the node at index $k - 1$ first.
That walk is $O(k)$, not $O(1)$. The list ADT's primitive is
\texttt{InsertAfter(curNode, newNode)} -- it assumes you already
have a pointer to the right spot. How you got that pointer (cached
from before, returned by a previous operation, walked to find)
isn't the ADT's concern.

This is why singly linked lists are often described as supporting
``$O(1)$ insertion \emph{given a pointer}'' -- the pointer is
assumed to be available.
\end{ideabox}

\subsection{Mapping to \texttt{std::list}}

\begin{examplebox}[The STL version]
\begin{lstlisting}[basicstyle=\ttfamily\small,frame=single,language=C++]
#include <list>
std::list<int> l = {7, 9, 5};

auto it = std::next(l.begin(), 1);   // iterator to the "9"
l.insert(it, 42);                    // inserts *before* it: 7, 42, 9, 5

auto it2 = std::next(l.begin(), 1);  // iterator to the "42"
l.insert(std::next(it2), 100);       // inserts before "9": 7, 42, 100, 9, 5
\end{lstlisting}
\texttt{std::list::insert(pos, x)} inserts \emph{before}
\texttt{pos}. To insert \emph{after} a position, pass
\texttt{std::next(pos)}. The STL picks ``before'' because that's
what \texttt{pos = end()} naturally means -- insert at the end.
Under the hood it's the same two-pointer splice, just with a
doubly linked list.
\end{examplebox}

\begin{notebox}[Why the pseudocode picks ``after'' instead of ``before'']
On a singly linked list, inserting \emph{before} a node is
surprisingly annoying -- you need a pointer to the node \emph{that
points to} the target, i.e.\ its predecessor. The list doesn't
provide that; you'd have to walk from \texttt{head} to find it,
making the operation $O(n)$. Inserting \emph{after} only needs the
target itself. That's why the SLL ADT exposes
\texttt{InsertAfter}, not \texttt{InsertBefore}. Doubly linked
lists, with back-pointers, get both for free.
\end{notebox}

\subsection{Cost}

\begin{center}
\footnotesize
\renewcommand{\arraystretch}{1.3}
\begin{tabular}{@{}lll@{}}
\hline
\textbf{Operation}    & \textbf{Singly linked list} & \textbf{Why} \\
\hline
\texttt{insertAfter(curNode, \ldots)} & $O(1)$      & Two pointer writes; no scan \\
Find \texttt{curNode} by value or index & $O(n)$    & Walk from head \\
\textbf{Insert by index $k$}            & $O(k) + O(1) = O(n)$ & Walk first, then splice \\
\hline
\end{tabular}
\end{center}

\begin{warnbox}[The $O(1)$ claim is only about the splice, not the find]
Everyone who says ``linked lists have $O(1)$ insertion'' is
technically right about the splice, but the \emph{find} is
$O(n)$. Unless you're iterating the list and splicing as you go, or
you already have a cached node pointer, the practical cost of
``insert an item at position $k$'' is linear. This is the biggest
source of confusion about linked-list performance.
\end{warnbox}

\section{4.4 Singly Linked Lists: Remove}

Removal in a singly linked list has the same shape of problem as
insertion: you need to update a \texttt{next} pointer that lives in
the \emph{predecessor} of the node you want to delete -- and a
singly linked list doesn't make that predecessor easy to reach. The
workaround is the same one used by \texttt{insertAfter}: name the
operation \textbf{RemoveAfter}, and make the caller pass the
predecessor pointer.

\begin{defnbox}[RemoveAfter]
\textbf{RemoveAfter(list, curNode)} removes the node
\emph{immediately after} \texttt{curNode}. Special case: if
\texttt{curNode} is \texttt{null}, the function removes the
\emph{head} (the first node has no predecessor inside the list
itself).
\end{defnbox}

\subsection{Why predecessor-based removal?}

\begin{ideabox}[The asymmetry that makes this awkward]
To remove node $X$ from a linked list you have to rewrite
\texttt{prev$\to$next} to skip over $X$ -- but in a singly linked
list, \texttt{prev} isn't available from $X$ itself; there's no back
pointer. The list has three realistic options:
\begin{enumerate}[leftmargin=*]
  \item \textbf{Walk from head} to find \texttt{prev}. $O(n)$.
  \item \textbf{Make the caller hand us \texttt{prev}} and expose
        \texttt{RemoveAfter(prev)}. $O(1)$.
  \item \textbf{Switch to a doubly linked list} where each node
        knows its predecessor. $O(1)$, but more memory.
\end{enumerate}
The textbook's ADT picks option 2. Option 3 is what we'll build in
4.5.
\end{ideabox}

\subsection{The three cases}

Two variables carry the state of the splice:
\begin{itemize}[leftmargin=*]
  \item \texttt{curNode} -- the node \emph{before} the one being
        removed (or \texttt{null} if we're removing the head).
  \item \texttt{sucNode} -- the node \emph{after} the one being
        removed (i.e., the new successor that the splice will
        connect to).
\end{itemize}

\begin{center}
\footnotesize
\renewcommand{\arraystretch}{1.3}
\begin{tabular}{@{}p{4cm}p{8cm}@{}}
\hline
\textbf{Case} & \textbf{What gets updated} \\
\hline
Remove the head (\texttt{curNode} null)  & \texttt{head} moves to successor; if that's null, also null out \texttt{tail}. \\
Remove an interior node                   & \texttt{curNode$\to$next} moves to successor. No endpoint changes. \\
Remove the tail (\texttt{curNode$\to$next$\to$next} null) & \texttt{curNode$\to$next} becomes null; also move \texttt{tail} to \texttt{curNode}. \\
\hline
\end{tabular}
\end{center}

\subsection{The algorithm}

\begin{examplebox}[\texttt{removeAfter}]
\begin{lstlisting}[basicstyle=\ttfamily\small,frame=single,language=C++]
void removeAfter(LinkedList& list, Node* curNode) {
    // Case A: remove the head
    if (curNode == nullptr && list.head != nullptr) {
        Node* toFree = list.head;
        Node* sucNode = list.head->next;
        list.head = sucNode;
        if (sucNode == nullptr) {        // list is now empty
            list.tail = nullptr;
        }
        delete toFree;
        return;
    }

    // Case B / C: remove the node after curNode (if one exists)
    if (curNode != nullptr && curNode->next != nullptr) {
        Node* toFree = curNode->next;
        Node* sucNode = curNode->next->next;
        curNode->next = sucNode;
        if (sucNode == nullptr) {        // we just removed the tail
            list.tail = curNode;
        }
        delete toFree;
    }
}
\end{lstlisting}
The textbook's pseudocode omits the \texttt{delete}. In real C++
you \emph{must} free the removed node, or you leak memory. Python,
Java, and JavaScript's garbage collectors would do it for you;
C++ makes you explicit.
\end{examplebox}

\begin{warnbox}[Forgetting \texttt{delete} is a memory leak]
Every \texttt{new} needs a matching \texttt{delete} unless a smart
pointer is handling the ownership. Hand-rolled linked lists are one
of the easiest places to leak memory in C++ because the node you
unlink is no longer reachable from \texttt{head}, so nothing else
will ever free it. \textbf{Cache the node before unlinking,
reassign the pointers, then \texttt{delete} the saved node.}
\end{warnbox}

\begin{warnbox}[Do not \texttt{delete} before you finish reading]
\texttt{Node* sucNode = curNode$\to$next$\to$next;} reads through
\texttt{curNode$\to$next}. If you \texttt{delete curNode$\to$next}
first, the subsequent dereference is undefined behavior. Order of
operations:
\begin{enumerate}[leftmargin=*]
  \item Cache the node to be deleted.
  \item Cache any pointers you'll need \emph{from} it.
  \item Rewrite the surrounding pointers.
  \item \texttt{delete} the cached node.
\end{enumerate}
This ``read-then-rewrite-then-delete'' rhythm is the heart of
safe pointer surgery.
\end{warnbox}

\subsection*{Picture of the interior splice}

\begin{center}
\footnotesize
Before:\\
\texttt{... $\to$ [curNode | $\bullet$] $\to$ [victim | $\bullet$] $\to$ [sucNode | ...] $\to$ ...}
\\[0.5em]
After \texttt{curNode$\to$next = sucNode}:\\
\texttt{... $\to$ [curNode | $\bullet$] $\to$ [sucNode | ...] $\to$ ...}
\\[0.5em]
The \texttt{victim} node is unreachable -- \texttt{delete} it.
\end{center}

\subsection{Cost}

\begin{center}
\footnotesize
\renewcommand{\arraystretch}{1.3}
\begin{tabular}{@{}lll@{}}
\hline
\textbf{Operation}                & \textbf{Singly linked list} & \textbf{Why} \\
\hline
\texttt{removeAfter(curNode)}     & $O(1)$                      & Two pointer writes; no scan \\
Remove by value (find + remove)   & $O(n)$                      & Linear walk to find the predecessor \\
Remove at index $k$               & $O(k) = O(n)$               & Walk to position, then splice \\
\hline
\end{tabular}
\end{center}

\subsection{Why we don't have \texttt{Remove(node)} directly}

\begin{ideabox}[The ``remove given just the victim'' trick]
There's a famous trick interview question: ``Given only a pointer
to the node you want to delete in a singly linked list (not the
head), delete it without a predecessor pointer.''

The trick: \emph{copy the next node's data into the victim and
delete the next node instead}. Works for interior nodes; fails for
the tail (there's no successor to pluck data from).

Real libraries avoid this contortion because it has subtle issues:
dangling references held by callers, data-copy cost if values are
large, ordering effects on equal-valued neighbors. \emph{Doubly
linked lists} solve the problem cleanly -- they give you the
predecessor pointer for free.
\end{ideabox}

\subsection{All five SLL operations together}

Now that we have append / prepend / insert-after / remove-after,
we can fill in the canonical cost table for a singly linked list
with both head and tail pointers:

\begin{center}
\footnotesize
\renewcommand{\arraystretch}{1.3}
\begin{tabular}{@{}lll@{}}
\hline
\textbf{Operation}   & \textbf{Cost}      & \textbf{Assumption} \\
\hline
\texttt{append(x)}        & $O(1)$         & Tail pointer cached \\
\texttt{prepend(x)}       & $O(1)$         & -- \\
\texttt{insertAfter(p, x)} & $O(1)$        & Caller has pointer \texttt{p} \\
\texttt{removeAfter(p)}   & $O(1)$         & Caller has pointer \texttt{p} \\
\texttt{search(x)}        & $O(n)$         & Linear scan \\
\texttt{at(i)}            & $O(i)$         & Walk from head \\
\texttt{remove(x)} by value & $O(n)$       & Search + splice \\
\hline
\end{tabular}
\end{center}

\begin{notebox}[\texttt{std::forward\_list} is this structure in the STL]
\texttt{std::forward\_list} is the C++ singly linked list. Its
iterators support \texttt{insert\_after} and \texttt{erase\_after}
for exactly the reason described here -- those are the $O(1)$
operations the structure can actually give you. It deliberately
omits \texttt{size()}, \texttt{push\_back}, and bidirectional
iteration because giving them would either force $O(n)$ costs or
require back-pointers (which would make it \texttt{std::list}).
\end{notebox}

\begin{warnbox}[Linked lists rarely win in modern C++]
Despite the clean $O(1)$ insertion/deletion, singly linked lists
lose to \texttt{std::vector} on most real workloads -- even ones
that intuitively seem to favor linked structure. The reasons are
cache locality (a vector is one contiguous block; linked nodes are
scattered) and per-node overhead (one pointer per item in an SLL,
two in a DLL, zero in a vector). Use a linked list when you
\emph{need} stable pointers into the middle of the list or when you
splice big ranges; otherwise reach for a vector.
\end{warnbox}

\section{4.5 Linked List Search}

Search on a linked list is exactly what you'd expect: start at the
head, walk node-by-node following \texttt{next} pointers, stop when
you find a match or hit \texttt{null}. This is \emph{linear search}
from Chapter 3, applied to a linked-list structure. It's worth its
own section because linked-list search exposes a recurring pattern
-- the \emph{walker idiom} -- that will be the scaffolding for
traversal, reversal, length computation, and many other linked-list
algorithms.

\begin{defnbox}[Linked-list search]
Given a list and a \textbf{key}, search walks from the head and
returns the first node whose \texttt{data} equals the key. If no
node matches, it returns \texttt{nullptr}.
\end{defnbox}

\subsection{The algorithm}

\begin{examplebox}[\texttt{search}]
\begin{lstlisting}[basicstyle=\ttfamily\small,frame=single,language=C++]
Node* search(const LinkedList& list, int key) {
    Node* curNode = list.head;
    while (curNode != nullptr) {
        if (curNode->data == key) {
            return curNode;
        }
        curNode = curNode->next;
    }
    return nullptr;
}
\end{lstlisting}
Three things to notice: the loop advances through
\texttt{curNode = curNode$\to$next}; the termination check is on
the pointer, not the data; and we return the \emph{pointer to the
node}, not its index.
\end{examplebox}

\begin{ideabox}[The walker idiom]
The structure \texttt{Node* curNode = list.head; while
(curNode$\neq$nullptr) \{ \ldots; curNode = curNode$\to$next; \}} is
the linked-list equivalent of \texttt{for (i = 0; i < n; ++i)} on
an array. Almost every linked-list algorithm in this chapter --
print, length, reverse, remove-by-value, insert-sorted -- is built
on this idiom. Memorize it.
\end{ideabox}

\subsection{Why return a pointer instead of an index}

\begin{warnbox}[Indices aren't the right return type for linked lists]
On a vector, \texttt{find} returns an index. On a linked list,
returning an index would be almost useless -- the caller would have
to walk from head \emph{again} to reach that position. Returning a
\emph{pointer to the node} gives the caller immediate access to the
data and to its neighbors for subsequent \texttt{InsertAfter} or
\texttt{RemoveAfter} calls. The ADT is designed around passing
these node pointers around.

The STL mirrors this: \texttt{std::find} on a vector returns an
iterator that's effectively a pointer into the array;
\texttt{std::find} on a list returns an iterator that wraps a node
pointer. Same \emph{idea}, different underlying representation --
which is exactly what an ADT/iterator abstraction is for.
\end{warnbox}

\subsection{Cost}

\begin{ideabox}[$O(n)$, always]
Linked-list search has to examine each node in turn until it finds
the key or exhausts the list. No structural shortcut is available
-- we can't binary-search a linked list, because reaching the
middle is $O(n/2) = O(n)$ on its own. Sorted or unsorted, the
linked-list search runtime is $\Theta(n)$ worst case.
\end{ideabox}

\begin{center}
\footnotesize
\renewcommand{\arraystretch}{1.25}
\begin{tabular}{@{}lll@{}}
\hline
\textbf{Case}    & \textbf{Comparisons} & \textbf{Runtime} \\
\hline
Key at head       & 1                   & $O(1)$ \\
Key at position $k$ & $k + 1$           & $O(k)$ \\
Key absent         & $n$                & $O(n)$ \\
\hline
\end{tabular}
\end{center}

\begin{warnbox}[Sorting a linked list doesn't help search]
On an array, sorting unlocks binary search: $O(\log n)$. On a
linked list, it does not. Even with a sorted linked list, you
can't jump to the middle element in constant time. The best you
can do is stop early when a value exceeds the key (``early exit''
on sorted data):
\begin{lstlisting}[basicstyle=\ttfamily\small,frame=single,language=C++]
while (curNode != nullptr && curNode->data < key) {
    curNode = curNode->next;
}
return (curNode && curNode->data == key) ? curNode : nullptr;
\end{lstlisting}
Helpful on hits near the front; still $\Theta(n)$ worst case.
\end{warnbox}

\subsection{Linked-list search is hostile to CPU caches}

\begin{ideabox}[Why array search beats linked search in practice]
Asymptotically, linear scan on a vector and linear search on a
linked list are both $\Theta(n)$. On real hardware they differ by
$10\times$ or more. The reason: a vector stores values in a single
contiguous block of memory, so the CPU's prefetcher can pull ahead
and cache the next dozens of elements before you ask for them. A
linked list's nodes can be scattered anywhere in the heap, and
following \texttt{next} pointers is a series of dependent loads
that the prefetcher cannot anticipate. Every \texttt{curNode =
curNode$\to$next} potentially stalls waiting on main memory.

This is the top reason modern code prefers \texttt{std::vector}
even when textbook analysis suggests \texttt{std::list} would work.
\end{ideabox}

\subsection{Two common variations}

\begin{examplebox}[Return a pointer \emph{and} its predecessor]
Sometimes a caller wants both the found node and the one before it
-- so they can immediately \texttt{RemoveAfter}. The idiom tracks
\texttt{prev} alongside \texttt{curNode}:
\begin{lstlisting}[basicstyle=\ttfamily\small,frame=single,language=C++]
std::pair<Node*, Node*> searchWithPrev(const LinkedList& list, int key) {
    Node* prev = nullptr;
    Node* cur  = list.head;
    while (cur != nullptr) {
        if (cur->data == key) return {cur, prev};
        prev = cur;
        cur  = cur->next;
    }
    return {nullptr, nullptr};
}
\end{lstlisting}
The \texttt{prev} pointer is the ``trailing pointer'' trick. It's
the closest a singly linked list gets to a doubly linked list's
back pointer, synthesized fresh on each walk.
\end{examplebox}

\begin{examplebox}[Recursive search]
\begin{lstlisting}[basicstyle=\ttfamily\small,frame=single,language=C++]
Node* searchRec(Node* cur, int key) {
    if (cur == nullptr)        return nullptr;
    if (cur->data == key)      return cur;
    return searchRec(cur->next, key);
}
\end{lstlisting}
Same algorithm, recursive shape. Intellectually clean, but each
call adds a stack frame -- $O(n)$ recursion depth. The iterative
version uses $O(1)$ stack. Prefer iterative for long lists;
recursive is just for illustration or when the language
(e.g.\ functional languages with tail-call optimization) supports
it cheaply.
\end{examplebox}

\begin{warnbox}[Linked lists and recursion don't mix well in C++]
A 100{,}000-node list processed with recursive \texttt{search} will
happily blow your stack. Tail-call elimination isn't guaranteed in
C++, and even when compilers do it, they don't always. Keep the
iterative version as your default for any traversal that could hit
long lists.
\end{warnbox}

\subsection{What search gives you}

\begin{ideabox}[The building block]
\texttt{search} is the first operation that doesn't already have a
cached pointer handed to it. Most higher-level operations you'd
write -- ``remove all nodes with value $x$,'' ``insert $y$ after
the first $x$,'' ``count how many $x$'s are in the list'' -- are
variations of search with different actions on match. The walker
idiom plus a small change at the match site gives you almost
everything.
\end{ideabox}

\section{4.6 Doubly Linked Lists}

The pain of a singly linked list lives at the \texttt{prev}
pointer it doesn't have: no cheap reverse traversal, no removal
without a predecessor argument, no insertion ``before'' a known
node. The \textbf{doubly linked list} spends one extra pointer per
node to erase all of these problems. You get bidirectional
traversal, $O(1)$ removal given the victim itself, and a cleaner
(if wordier) splice operation.

\begin{defnbox}[Doubly linked list]
A \textbf{doubly linked list (DLL)} is a linked list where each
node has two pointers: \texttt{next} to its successor and
\texttt{prev} to its predecessor. The list still holds \texttt{head}
and \texttt{tail} pointers; the head node's \texttt{prev} is
\texttt{null} and the tail node's \texttt{next} is \texttt{null}.
\end{defnbox}

\begin{ideabox}[The picture]
\texttt{null $\leftarrow$ [$\bullet$ | 7 | $\bullet$] $\leftrightarrow$ [$\bullet$ | 9 | $\bullet$] $\leftrightarrow$ [$\bullet$ | 5 | $\bullet$] $\to$ null}

Each cell has \texttt{prev} on the left and \texttt{next} on the
right. The bidirectional arrows capture that every link is
traversable both ways.
\end{ideabox}

\subsection{The node type in C++}

\begin{examplebox}[A DLL node]
\begin{lstlisting}[basicstyle=\ttfamily\small,frame=single,language=C++]
struct DNode {
    int    data;
    DNode* next;
    DNode* prev;

    DNode(int d) : data(d), next(nullptr), prev(nullptr) {}
};

struct DLinkedList {
    DNode* head = nullptr;
    DNode* tail = nullptr;
};
\end{lstlisting}
One extra pointer per node compared to an SLL. On a 64-bit system
that's 8 extra bytes per item; for 32-bit \texttt{int} data plus
two 64-bit pointers, a DLL node is already 20 bytes of pointers
and padding wrapping 4 bytes of data -- a 5$\times$ memory
overhead versus a \texttt{std::vector<int>}. Bring this up when
someone claims linked lists are ``cheap.''
\end{examplebox}

\subsection{Invariants}

\begin{ideabox}[What a well-formed DLL always satisfies]
\begin{enumerate}[leftmargin=*]
  \item If \texttt{node$\to$next} is not null,
        \texttt{node$\to$next$\to$prev} equals \texttt{node}.
  \item If \texttt{node$\to$prev} is not null,
        \texttt{node$\to$prev$\to$next} equals \texttt{node}.
  \item \texttt{head$\to$prev == nullptr}, \texttt{tail$\to$next == nullptr}.
  \item \texttt{head == nullptr} iff \texttt{tail == nullptr}.
\end{enumerate}
Every operation below has to restore all four invariants. A DLL
bug is almost always a forgotten back-pointer update.
\end{ideabox}

\subsection{Append}

\begin{examplebox}[\texttt{append} on a DLL]
\begin{lstlisting}[basicstyle=\ttfamily\small,frame=single,language=C++]
void append(DLinkedList& list, DNode* newNode) {
    if (list.head == nullptr) {          // empty list
        list.head = newNode;
        list.tail = newNode;
    } else {                              // non-empty list
        list.tail->next = newNode;        // stitch forward
        newNode->prev   = list.tail;      // stitch backward
        list.tail       = newNode;        // advance tail
    }
}
\end{lstlisting}
Compared to the SLL version, we do \emph{one extra write}: set
\texttt{newNode$\to$prev = list.tail} before moving the tail. Still
$O(1)$.
\end{examplebox}

\begin{warnbox}[Every next-write deserves a twin prev-write]
The rule of thumb when implementing any DLL operation: for every
\texttt{a$\to$next = b} you write, ask ``did I also write
\texttt{b$\to$prev = a}?'' Skip that check and you'll silently break
invariant 1. Many DLL bugs only show up when iterating the list
\emph{backwards} -- which can be long after the broken write.
\end{warnbox}

\subsection{Prepend}

Symmetric to append. Every \texttt{tail} becomes \texttt{head};
every \texttt{next} swaps with \texttt{prev}.

\begin{examplebox}[\texttt{prepend} on a DLL]
\begin{lstlisting}[basicstyle=\ttfamily\small,frame=single,language=C++]
void prepend(DLinkedList& list, DNode* newNode) {
    if (list.head == nullptr) {          // empty list
        list.head = newNode;
        list.tail = newNode;
    } else {                              // non-empty list
        newNode->next  = list.head;       // forward link in
        list.head->prev = newNode;        // backward link in
        list.head       = newNode;        // advance head
    }
}
\end{lstlisting}
\end{examplebox}

\subsection{What the back-pointer unlocks}

\begin{ideabox}[Why the extra pointer is worth it]
\begin{itemize}[leftmargin=*]
  \item \textbf{Reverse iteration.} Walk from \texttt{tail}
        following \texttt{prev}. SLL can't do this without
        reversing the list or recursing.
  \item \textbf{\texttt{Remove(node)} with just the victim.} Given a
        pointer to any node, you can unlink it without having the
        predecessor cached -- you read \texttt{node$\to$prev} to
        find it. SLL requires \texttt{removeAfter(prev)}.
  \item \textbf{\texttt{InsertBefore(node, x)}.} Same reason:
        \texttt{node$\to$prev} gives you everything needed to
        splice before.
  \item \textbf{\texttt{PrintReverse}, \texttt{Sort} with
        bidirectional access, iterators in both directions.}
\end{itemize}
Every common list operation is either equivalent in cost to SLL or
strictly easier on a DLL.
\end{ideabox}

\subsection{When the extra pointer is \emph{not} worth it}

\begin{warnbox}[Two places SLL still wins]
\begin{itemize}[leftmargin=*]
  \item \textbf{Memory-sensitive scenarios.} One pointer per node
        instead of two is 50\% less pointer overhead. If you're
        storing millions of tiny items, this actually matters.
        \texttt{std::forward\_list} exists precisely for this.
  \item \textbf{Intrusive linked lists in kernels or embedded
        systems.} Fewer pointers mean smaller per-object headers.
\end{itemize}
Outside those niches, the DLL wins on ergonomics. \texttt{std::list}
is a DLL, not an SLL, for this reason.
\end{warnbox}

\subsection{Cost of the endpoint operations}

\begin{center}
\footnotesize
\renewcommand{\arraystretch}{1.3}
\begin{tabular}{@{}llll@{}}
\hline
\textbf{Operation}    & \textbf{SLL} & \textbf{DLL} & \textbf{Notes} \\
\hline
\texttt{append}       & $O(1)$       & $O(1)$        & One extra write on DLL \\
\texttt{prepend}      & $O(1)$       & $O(1)$        & One extra write on DLL \\
\texttt{insertAfter(n)} & $O(1)$    & $O(1)$        & DLL also has \texttt{insertBefore} \\
\texttt{remove(n)}   & needs prev, so $O(n)$ to find it, then $O(1)$ & $O(1)$ given node & DLL wins outright \\
\hline
\end{tabular}
\end{center}

\subsection{\texttt{std::list} is a circular DLL}

\begin{notebox}[How the STL implements DLL]
\texttt{std::list} uses a circular doubly linked list with a
\emph{sentinel node}. A dedicated sentinel (not real data) sits
between the tail and the head; the list is never empty in the
linked-structure sense -- even an empty list has the sentinel
pointing at itself. This trick eliminates all the ``is the list
empty?'' special cases and gives \texttt{end()} a stable iterator
to return, but you won't write it by hand in an intro course. The
plain head/tail/null-terminator DLL here is the standard pedagogy.
\end{notebox}

\subsection{Where this is going}

\begin{ideabox}[The rest of the DLL arc]
Appending and prepending are the warm-up. The interesting DLL
operations come next:
\begin{itemize}[leftmargin=*]
  \item \textbf{InsertAfter} (4.7 if the textbook follows the
        pattern): same as SLL plus one back-pointer update per side.
  \item \textbf{Remove(node)}: the big payoff -- no
        \texttt{prev}-arg needed.
  \item \textbf{PrintReverse}: trivial on DLL, annoying on SLL.
\end{itemize}
Once you've seen DLL's remove, the trade-off between SLL and DLL
should be obvious: you pay one extra pointer per node for a much
cleaner API. For 99\% of production code, you take the deal.
\end{ideabox}

\begin{warnbox}[C++ memory hygiene reminder]
The same \texttt{delete} discipline from the SLL remove section
applies here: save the victim pointer, unlink all four
neighboring pointers (two \texttt{next}, two \texttt{prev}),
\texttt{delete} the victim. Smart pointers (\texttt{std::unique\_ptr})
complicate DLL ownership because both \texttt{next} and \texttt{prev}
``own'' each other -- don't naively wrap pointers in
\texttt{unique\_ptr} on both sides or you'll get cyclic ownership.
Typical practice: \texttt{prev} is a raw pointer; \texttt{next}
is a \texttt{unique\_ptr}. Or, skip this and just use
\texttt{std::list}.
\end{warnbox}

\section{4.7 Doubly Linked Lists: Insert-After}

Insert-after on a doubly linked list is the same algorithm as the
SLL version from 4.3, with one extra pointer write per affected
side to keep the \texttt{prev} links in sync. The interior case --
where all four invariants meet -- requires \textbf{four} pointer
updates instead of the SLL's two. Each of the four has to happen in
the right order, and every pattern you pick here will be the
template for the rest of the DLL operations.

\begin{defnbox}[InsertAfter on a DLL]
\textbf{InsertAfter(list, curNode, newNode)} inserts
\texttt{newNode} immediately after \texttt{curNode}. As in the SLL
version, the three cases are ``empty list,'' ``\texttt{curNode} is
the tail,'' and ``interior.'' Each case preserves the DLL's four
invariants (both directions linked, head/tail nulls, empty-list
agreement).
\end{defnbox}

\subsection{The three cases}

\begin{center}
\footnotesize
\renewcommand{\arraystretch}{1.3}
\begin{tabular}{@{}p{3.8cm}p{3cm}p{5cm}@{}}
\hline
\textbf{Case} & \textbf{Pointer writes} & \textbf{Structures affected} \\
\hline
Empty list                      & 2 & \texttt{head}, \texttt{tail} \\
\texttt{curNode} is tail        & 3 & tail's \texttt{next}, new node's \texttt{prev}, \texttt{tail} \\
Interior                        & 4 & curNode's \texttt{next}, sucNode's \texttt{prev}, new node's \texttt{next} and \texttt{prev} \\
\hline
\end{tabular}
\end{center}

\subsection{The algorithm}

\begin{examplebox}[\texttt{insertAfter} on a DLL]
\begin{lstlisting}[basicstyle=\ttfamily\small,frame=single,language=C++]
void insertAfter(DLinkedList& list, DNode* curNode, DNode* newNode) {
    if (list.head == nullptr) {                      // Case A: empty
        list.head = newNode;
        list.tail = newNode;
    }
    else if (curNode == list.tail) {                  // Case B: after tail
        list.tail->next = newNode;
        newNode->prev   = list.tail;
        list.tail       = newNode;
    }
    else {                                             // Case C: interior
        DNode* sucNode = curNode->next;
        newNode->next  = sucNode;
        newNode->prev  = curNode;
        curNode->next  = newNode;
        sucNode->prev  = newNode;
    }
}
\end{lstlisting}
\end{examplebox}

\subsection*{Case C, step by step}

The interior case is four pointer writes. The order is
\emph{flexible}, but the intent must stay consistent: \textbf{wire
the new node first, then redirect the existing ones}. Writing the
old pointers first while the new node is still dangling breaks
invariants mid-flight and can get you into trouble on multi-threaded
or signal-interrupted code. Safe ordering:

\begin{center}
\footnotesize
\renewcommand{\arraystretch}{1.3}
\begin{tabular}{@{}rl@{}}
\hline
\textbf{Step} & \textbf{Write} \\
\hline
1 & \texttt{newNode$\to$next = curNode$\to$next}  (remember the successor first) \\
2 & \texttt{newNode$\to$prev = curNode}            (wire the predecessor) \\
3 & \texttt{curNode$\to$next = newNode}            (splice forward) \\
4 & \texttt{sucNode$\to$prev = newNode}            (splice backward) \\
\hline
\end{tabular}
\end{center}

\begin{warnbox}[Don't overwrite \texttt{curNode$\to$next} before caching it]
The \texttt{sucNode} local variable exists precisely to cache
\texttt{curNode$\to$next} \emph{before} we overwrite it in step 3.
If you skip the cache -- write step 3 first and then try to reach
\texttt{sucNode} via \texttt{curNode$\to$next} -- you've lost the
successor. Classic prerequisite for a ``why are half my nodes
missing'' debugging session. \textbf{Cache out before you
overwrite.}
\end{warnbox}

\subsection{Picture of the interior splice}

\begin{center}
\footnotesize
Before:\\
\texttt{... $\leftrightarrow$ [curNode] $\leftrightarrow$ [sucNode] $\leftrightarrow$ ...}
\\[0.5em]
After:\\
\texttt{... $\leftrightarrow$ [curNode] $\leftrightarrow$ [newNode] $\leftrightarrow$ [sucNode] $\leftrightarrow$ ...}
\end{center}

Four arrows redirected. Nothing else changes -- the head, the tail,
and any nodes further away are untouched.

\subsection{Why inserting \emph{before} is now trivial}

\begin{ideabox}[InsertBefore for free]
On a singly linked list, \texttt{InsertBefore(node, x)} requires
finding the predecessor -- a walk from \texttt{head}. On a doubly
linked list, \texttt{node$\to$prev} is right there.
\texttt{InsertBefore(node, x)} is \texttt{InsertAfter(node$\to$prev, x)}
with one special case for \texttt{node == head}. That's why
\texttt{std::list} and almost all modern container APIs use
``insert \emph{before} position'' as their primitive: on a DLL
it's natural.
\end{ideabox}

\begin{examplebox}[\texttt{insertBefore} on a DLL]
\begin{lstlisting}[basicstyle=\ttfamily\small,frame=single,language=C++]
void insertBefore(DLinkedList& list, DNode* curNode, DNode* newNode) {
    if (list.head == nullptr) {
        list.head = newNode;
        list.tail = newNode;
    }
    else if (curNode == list.head) {                   // before the head
        newNode->next   = list.head;
        list.head->prev = newNode;
        list.head       = newNode;
    }
    else {
        DNode* predNode = curNode->prev;
        newNode->next  = curNode;
        newNode->prev  = predNode;
        predNode->next = newNode;
        curNode->prev  = newNode;
    }
}
\end{lstlisting}
Symmetric to \texttt{insertAfter}, mirror-imaged. Same structure,
same four-writes-in-interior-case, same cache-predecessor-first
discipline.
\end{examplebox}

\subsection{Cost}

\begin{center}
\footnotesize
\renewcommand{\arraystretch}{1.3}
\begin{tabular}{@{}lll@{}}
\hline
\textbf{Operation}                   & \textbf{Cost} & \textbf{Assumption} \\
\hline
\texttt{insertAfter(curNode, x)}     & $O(1)$        & Caller has \texttt{curNode} \\
\texttt{insertBefore(curNode, x)}    & $O(1)$        & Caller has \texttt{curNode} \\
\texttt{insert at index $k$}          & $O(k)$        & Walk from head, then splice \\
\hline
\end{tabular}
\end{center}

\begin{warnbox}[``Four writes'' vs.\ SLL's ``two'' is not a real disadvantage]
A DLL insertion does 2$\times$ the pointer writes of an SLL
insertion, but both are constant-time. On modern hardware the extra
writes are lost in the noise (pipelined, same cache line). What
matters is that DLL inserts still don't scan or shift. You trade a
few nanoseconds per insert for a vastly more useful API -- cheap
reverse iteration, cheap removal from an arbitrary node, cheap
insertBefore.
\end{warnbox}

\subsection{In the STL}

\begin{examplebox}[\texttt{std::list::insert}]
\begin{lstlisting}[basicstyle=\ttfamily\small,frame=single,language=C++]
std::list<int> l = {7, 9, 5};
auto it = std::next(l.begin(), 1);  // iterator to "9"

l.insert(it, 42);                    // 7, 42, 9, 5  (insert before)
l.insert(std::next(it), 100);        // 7, 42, 9, 100, 5 (insert after "9")
\end{lstlisting}
\texttt{std::list::insert(pos, x)} inserts \emph{before}
\texttt{pos}. This is the only primitive the STL provides -- insert
after is written as \texttt{insert(std::next(pos), x)}. The
implementation under the hood is the four-pointer-write interior
case above, wrapped in a circular-DLL-with-sentinel that also
eliminates the ``empty list'' special case.
\end{examplebox}

\begin{ideabox}[Where this is going]
With insertAfter done, the remaining operation is \textbf{remove}.
The payoff of the DLL structure shows up there most clearly: with
\texttt{prev} pointers available, removal given just the victim
takes constant time and doesn't need a \texttt{removeAfter} trick.
Section 4.8 (or whatever number the textbook assigns it) will make
that concrete.
\end{ideabox}

\section{4.8 Doubly Linked Lists: Remove}

This is where the extra \texttt{prev} pointer pays its rent.
\textbf{On a DLL, removing a node takes $O(1)$ given nothing but the
node itself} -- no predecessor argument, no backward scan. The SLL
couldn't do this. The algorithm is also structurally cleaner:
instead of three scenario branches, we use four \emph{independent}
checks that each handle one of the pointers that might need
updating.

\begin{defnbox}[Remove on a DLL]
\textbf{Remove(list, curNode)} unlinks \texttt{curNode} from the
list and returns it (or frees it, depending on ownership). It
updates up to four pointers: the predecessor's \texttt{next}, the
successor's \texttt{prev}, and -- if \texttt{curNode} was at an
endpoint -- the list's \texttt{head} or \texttt{tail}.
\end{defnbox}

\subsection{Four independent checks instead of three cases}

The SLL's \texttt{RemoveAfter} had three cases. The DLL's
\texttt{Remove} is organized differently: four independent
``if this pointer exists, patch it'' checks. Any combination of
them fires depending on where the victim sits.

\begin{center}
\footnotesize
\renewcommand{\arraystretch}{1.3}
\begin{tabular}{@{}p{5.3cm}p{6.7cm}@{}}
\hline
\textbf{Condition} & \textbf{Action} \\
\hline
\texttt{sucNode $\neq$ null}    & \texttt{sucNode$\to$prev = predNode} \\
\texttt{predNode $\neq$ null}   & \texttt{predNode$\to$next = sucNode} \\
\texttt{curNode == list.head}   & \texttt{list.head = sucNode} \\
\texttt{curNode == list.tail}   & \texttt{list.tail = predNode} \\
\hline
\end{tabular}
\end{center}

\begin{ideabox}[Why four conditions cover all scenarios]
\begin{itemize}[leftmargin=*]
  \item Interior node: \texttt{sucNode} and \texttt{predNode} both
        exist; \texttt{curNode} is neither head nor tail. Only the
        first two conditions fire.
  \item Head removal, list has other nodes: \texttt{sucNode}
        exists, \texttt{predNode} is null, \texttt{curNode} is
        head. Conditions 1 and 3 fire.
  \item Tail removal, list has other nodes: \texttt{predNode}
        exists, \texttt{sucNode} is null, \texttt{curNode} is tail.
        Conditions 2 and 4 fire.
  \item Only node in the list: both \texttt{sucNode} and
        \texttt{predNode} are null; \texttt{curNode} is both head
        and tail. Only conditions 3 and 4 fire, setting both to
        \texttt{nullptr}. The list ends up empty.
\end{itemize}
Independence is elegant: each check restores one invariant, and
the whole algorithm is just running all four.
\end{ideabox}

\subsection{The algorithm}

\begin{examplebox}[\texttt{remove} on a DLL]
\begin{lstlisting}[basicstyle=\ttfamily\small,frame=single,language=C++]
void remove(DLinkedList& list, DNode* curNode) {
    DNode* sucNode  = curNode->next;
    DNode* predNode = curNode->prev;

    if (sucNode  != nullptr) sucNode->prev  = predNode;
    if (predNode != nullptr) predNode->next = sucNode;

    if (curNode == list.head) list.head = sucNode;
    if (curNode == list.tail) list.tail = predNode;

    delete curNode;
}
\end{lstlisting}
Six lines of logic, four conditional pointer writes, one
\texttt{delete}. No traversal. $O(1)$ in all cases.
\end{examplebox}

\begin{warnbox}[Read \texttt{curNode$\to$next} and \texttt{curNode$\to$prev} \emph{before} \texttt{delete}ing]
As in the SLL case: cache the neighbors out of \texttt{curNode}
first, \emph{then} free the node. Reading through a freed pointer
is undefined behavior. The order of operations in the code above is
load neighbors $\to$ patch them $\to$ fix head/tail $\to$
\texttt{delete}. Follow that sequence and you'll be safe.
\end{warnbox}

\subsection{Picture}

\begin{center}
\footnotesize
Before, for an interior removal:\\
\texttt{... $\leftrightarrow$ [predNode] $\leftrightarrow$ [curNode] $\leftrightarrow$ [sucNode] $\leftrightarrow$ ...}
\\[0.5em]
After:\\
\texttt{... $\leftrightarrow$ [predNode] $\leftrightarrow$ [sucNode] $\leftrightarrow$ ...}
\end{center}

The \texttt{curNode} is short-circuited by updating exactly two
pointers: \texttt{predNode$\to$next} and \texttt{sucNode$\to$prev}.
The victim is then unreachable from the list and can be safely
freed.

\subsection{The big comparison to SLL}

This is the place where the SLL/DLL trade-off lands most clearly.

\begin{center}
\footnotesize
\renewcommand{\arraystretch}{1.3}
\begin{tabular}{@{}p{4.3cm}p{3.5cm}p{3.5cm}@{}}
\hline
\textbf{Operation}             & \textbf{SLL}          & \textbf{DLL} \\
\hline
Remove, given node pointer     & $O(n)$ (scan for prev) or use awkward \texttt{RemoveAfter(prev)} & $O(1)$ from the node \\
Remove by value                 & $O(n)$ (walk to find) & $O(n)$ (walk to find) + $O(1)$ unlink \\
Remove during a walk           & $O(1)$ with trailing prev & $O(1)$ natively \\
Splice a node between two existing lists & painful; need prev in both places & trivial with the node's prev/next pointers \\
\hline
\end{tabular}
\end{center}

\begin{ideabox}[The ``remove during a walk'' idiom]
On a DLL you can walk forward and decide-to-remove as you go
without tracking a trailing pointer -- \texttt{curNode$\to$prev}
gives it to you. This is why \texttt{std::list::remove\_if} and
related algorithms are so clean on \texttt{std::list}: remove
wherever you like, the pointers are always available.
\end{ideabox}

\subsection{Cost}

\begin{center}
\footnotesize
\renewcommand{\arraystretch}{1.3}
\begin{tabular}{@{}lll@{}}
\hline
\textbf{Operation}          & \textbf{Cost} & \textbf{Assumption} \\
\hline
\texttt{remove(node)}       & $O(1)$        & Caller has a pointer to the node \\
Remove by value              & $O(n)$        & Linear search to find, then $O(1)$ unlink \\
Remove at index $k$          & $O(k)$        & Walk to position, then $O(1)$ unlink \\
Remove all of value $x$      & $O(n)$        & Single walk; remove matches in place \\
\hline
\end{tabular}
\end{center}

\subsection{Handling ownership in C++}

\begin{warnbox}[Who frees the node?]
If \texttt{remove} calls \texttt{delete curNode}, the caller must
not \texttt{delete} it again. But if you want the caller to take
ownership (say, to reinsert the node into a different list),
\texttt{remove} should return the pointer instead of deleting. A
clean API picks one -- mixing both leads to double-frees or leaks.
\texttt{std::list::erase} deletes; \texttt{std::list::splice}
moves. Pick the right primitive for the task.
\end{warnbox}

\begin{examplebox}[Return-the-node variant]
\begin{lstlisting}[basicstyle=\ttfamily\small,frame=single,language=C++]
DNode* unlink(DLinkedList& list, DNode* curNode) {
    DNode* sucNode  = curNode->next;
    DNode* predNode = curNode->prev;
    if (sucNode  != nullptr) sucNode->prev  = predNode;
    if (predNode != nullptr) predNode->next = sucNode;
    if (curNode == list.head) list.head = sucNode;
    if (curNode == list.tail) list.tail = predNode;
    curNode->next = curNode->prev = nullptr;   // disconnect for safety
    return curNode;                            // caller owns it now
}
\end{lstlisting}
The \texttt{curNode$\to$next = curNode$\to$prev = nullptr} line
matters: leaving stale pointers on a detached node invites bugs
where the caller accidentally walks from it back into the old list.
\end{examplebox}

\subsection{Where this leaves the list ADT}

\begin{ideabox}[The SLL/DLL decision, in one line]
\emph{Singly linked list} is the minimal implementation of the list
ADT. \emph{Doubly linked list} spends one extra pointer per node to
make every operation that touches a known node constant time and
every iteration bidirectional. For almost every real use case, take
the deal.
\end{ideabox}

\begin{notebox}[Where this is going]
Chapter 4's next arc moves from generic lists to \textbf{restricted}
list-based ADTs: stacks, queues, and deques. Each one is ``a list
with some operations removed so the remaining ones can be made
cheaper or more predictable.'' Stacks and queues are both usually
implemented as either a dynamic array or a (doubly) linked list;
the work you've done here is the foundation for both.
\end{notebox}

\section{4.9 Linked List Traversal}

Traversal is the loop that ties linked lists back to every other
container: ``do something to each element.'' It is also the cheapest
linked list operation to get right, because a linked list has
exactly one shape of walk through it --- \texttt{curNode =
curNode->next}, repeated until the pointer falls off the end.

\begin{defnbox}[Traversal]
An algorithm that visits every node in a list exactly once and
performs some operation on each. The ``operation'' is a parameter
of the traversal, not part of it: print, sum, check a predicate,
copy into another container --- the walk is the same.
\end{defnbox}

\subsection*{The walker idiom, one more time}

You have written this loop (or a close cousin) in every earlier
section. It deserves its own name because it is the single
building block under search, insert-after-value, remove-by-value,
size, clear, copy, and almost any other ``go look at everything''
operation.

\begin{examplebox}[Forward traversal, the loop everyone writes]
\begin{lstlisting}[language=C++,basicstyle=\ttfamily\small,frame=single]
// Apply `visit` to every node in a singly-
// or doubly-linked list, front to back.
void traverse(Node* head) {
    for (Node* cur = head; cur != nullptr; cur = cur->next) {
        visit(cur);
    }
}
\end{lstlisting}
\end{examplebox}

Three things make this loop airtight. First, the starting state is
the list handle (\texttt{head}), not a guess --- if the list is
empty, \texttt{head} is \texttt{nullptr} and the loop body never
runs. Second, the termination condition is \texttt{cur != nullptr},
which works for length 0, length 1, and length n with no special
cases. Third, the advance step is a single field read;
\texttt{next} is either a valid node or the sentinel
\texttt{nullptr} that the list's append/prepend/remove code is
contractually required to maintain.

\begin{warnbox}[Never dereference \texttt{cur} inside the condition]
A common rookie loop is \texttt{while (cur->next != nullptr)},
which walks \emph{between} pairs of nodes. That form is useful for
\emph{predecessor-finding} (e.g., ``find the node before the one
to delete''), but it is the \textbf{wrong} loop for a traversal.
Using it for traversal silently skips the tail node. Match the
loop shape to the question you are answering.
\end{warnbox}

\subsection*{Doubly-linked lists: reverse traversal for free}

Because every node in a DLL carries a \texttt{prev} pointer and the
list keeps a \texttt{tail} handle, you can walk the list backwards
at the same cost as forwards. This is one of the concrete payoffs
for the extra pointer per node discussed in section~4.6.

\begin{examplebox}[Reverse traversal on a DLL]
\begin{lstlisting}[language=C++,basicstyle=\ttfamily\small,frame=single]
void traverseReverse(DNode* tail) {
    for (DNode* cur = tail; cur != nullptr; cur = cur->prev) {
        visit(cur);
    }
}
\end{lstlisting}
\end{examplebox}

A singly-linked list has no efficient reverse traversal. The only
way to visit the nodes in reverse is to either reverse the list
first (O(n) destructive), push every node onto a stack during a
forward pass (O(n) space), or use recursion (which is push-onto-
a-stack by another name, using the call stack). All three cost
Theta(n) extra work or space on top of the walk itself. The DLL
does it in zero extra work because the backward links were already
built.

\subsection*{Cost}

\begin{center}
\begin{tabular}{|l|c|l|}
\hline
\textbf{Operation} & \textbf{Cost} & \textbf{Notes} \\
\hline
Forward traversal (SLL or DLL) & $\Theta(n)$ & Must visit each node \\
Reverse traversal (DLL)        & $\Theta(n)$ & Symmetric to forward \\
Reverse traversal (SLL)        & $\Theta(n)$ & Plus $\Theta(n)$ space \\
\hline
\end{tabular}
\end{center}

The big-Theta is non-negotiable: to touch every element you pay
n operations. The interesting cost is \emph{per-step}, which is
O(1) for both SLL and DLL because each step is one pointer
indirection. This is what makes the linked list traversal the
same asymptotic cost as an array's \texttt{for (int i = 0; i < n;
++i)} even though the underlying machine work is quite different.

\begin{notebox}[Array traversal is faster in practice]
A linked-list traversal is $\Theta(n)$, but so is an array's
\texttt{for} loop --- and on real hardware the array wins by a
large constant factor. Array elements are contiguous, so the CPU
prefetches the next cache line while the current one is being
processed; linked-list nodes are scattered across the heap, so
each \texttt{cur = cur->next} is likely to miss the cache and
stall. Big-O says they're equal. The wall clock says otherwise.
This is one reason \texttt{std::vector} is the default C++
container even when the textbook problem ``wants'' a list.
\end{notebox}

\subsection*{Traversal as the hidden primitive}

Once you have the walker loop, almost every other ``touch
everything'' operation is a specialization of it.

\begin{center}
\begin{tabular}{|l|l|}
\hline
\textbf{Operation} & \textbf{What the walker body does} \\
\hline
\texttt{size()}         & Increment a counter \\
\texttt{print()}        & Print \texttt{cur->data} \\
\texttt{search(x)}      & Compare, return if found \\
\texttt{sum()}          & Accumulate \texttt{cur->data} \\
\texttt{clear()}        & \texttt{delete prev}, advance \\
\texttt{copy()}         & Append \texttt{cur->data} to new list \\
\texttt{contains(x)}    & Compare, return \texttt{bool} \\
\hline
\end{tabular}
\end{center}

The reason linked list code looks repetitive across the chapter
is that it \emph{is} repetitive: every one of these operations
is the same loop with a different body. If you find yourself
writing the traversal scaffold for the fifth time in one class,
that is the signal to factor out a higher-order helper that takes
a callable --- which is exactly what the STL does in
\texttt{std::for\_each}, \texttt{std::find\_if}, and
\texttt{std::accumulate}.

\begin{examplebox}[STL: traversal without writing the loop]
\begin{lstlisting}[language=C++,basicstyle=\ttfamily\small,frame=single]
#include <list>
#include <algorithm>
#include <numeric>
#include <iostream>

std::list<int> xs = {3, 1, 4, 1, 5, 9};

// print each
std::for_each(xs.begin(), xs.end(),
              [](int v) { std::cout << v << ' '; });

// find first element > 4
auto it = std::find_if(xs.begin(), xs.end(),
                       [](int v) { return v > 4; });

// sum
int total = std::accumulate(xs.begin(), xs.end(), 0);
\end{lstlisting}
\end{examplebox}

\texttt{std::list} exposes iterators (\texttt{begin()},
\texttt{end()}, \texttt{rbegin()}, \texttt{rend()}) that are the
library's name for ``the current walker.'' Advancing an iterator
with \texttt{++it} is exactly the \texttt{cur = cur->next} step;
the iterator just hides the pointer. Once you understand the raw
loop, the STL version stops looking magical.

\begin{warnbox}[Do not modify the list \emph{during} a traversal]
If the visit step can delete or re-link the current node, the
naive loop is broken --- you will read \texttt{cur->next} from a
node you just \texttt{delete}'d. The fix is the cache-before-
overwrite pattern from the remove sections: grab
\texttt{Node* nxt = cur->next;} before doing anything destructive,
then advance to \texttt{nxt} at the end of the loop body. The
STL's erase-returns-the-next-iterator idiom
(\texttt{it = xs.erase(it)}) is the same pattern with iterator
dress.
\end{warnbox}

\begin{ideabox}[Why this is the simplest section in the chapter]
Every earlier linked-list operation had edge cases: empty list,
single node, head, tail, interior, value-not-present. Traversal
has \textbf{none} of those. The loop works for length 0 because
\texttt{head == nullptr}. It works for length 1 because the
single node's \texttt{next} is \texttt{nullptr}. It works for
length n for the same reason. That simplicity is why every other
linked-list algorithm is built on top of it --- get the walker
right once and reuse it forever.
\end{ideabox}

\section{4.10 Sorting Linked Lists}

Sorting is the first algorithm in this chapter that isn't just a
walker variant. It's also where the cost difference between arrays
and linked lists stops being theoretical --- many classic array
sorts straight up don't work on a linked list, and the ones that
do have a different implementation on each list flavor. Insertion
sort was first analyzed on arrays in \textsection 3.15; here we
adapt it to a linked-list traversal.

\begin{ideabox}[The core question]
Which sorting algorithms transplant cleanly to a linked list, and
which need random access so badly they fall apart without it? The
answer splits neatly: algorithms that walk the data once or twice
(insertion, merge) transfer; algorithms that jump around by index
(quick, heap, shell) do not.
\end{ideabox}

\subsection*{Insertion sort on a DLL: the easy case}

Insertion sort on a doubly-linked list is almost identical to the
array version. The outer loop visits each node starting from the
second; the inner loop walks \texttt{prev} pointers back through
the sorted prefix until it finds the right slot, then splices the
current node there.

\begin{examplebox}[DLL insertion sort]
\begin{lstlisting}[language=C++,basicstyle=\ttfamily\small,frame=single]
void insertionSort(DList& L) {
    DNode* cur = L.head ? L.head->next : nullptr;
    while (cur) {
        DNode* nxt = cur->next;   // cache before mutation
        DNode* s   = cur->prev;
        while (s && s->data > cur->data) s = s->prev;

        L.remove(cur);
        if (!s) L.prepend(cur);
        else    L.insertAfter(s, cur);

        cur = nxt;
    }
}
\end{lstlisting}
\end{examplebox}

The cache-\texttt{nxt}-before-mutating pattern is the same one you
saw in remove-during-traversal (section 4.9). You can't advance
via \texttt{cur->next} after \texttt{L.remove(cur)} splices
\texttt{cur} out --- its \texttt{next} may now point somewhere
else, or may be stale.

\begin{warnbox}[SLL insertion sort is structurally different]
An SLL has no \texttt{prev} link, so you \emph{cannot} walk
backward from \texttt{cur} through the sorted prefix. The SLL
version flips the search: walk from \texttt{head} forward until
you find the insertion point, then splice. Same O(n$^2$)
asymptotic, different loop structure. Don't assume a doubly-linked
algorithm ports to singly-linked; check that every pointer read
you need actually exists.
\end{warnbox}

\begin{examplebox}[SLL insertion sort, forward-search variant]
\begin{lstlisting}[language=C++,basicstyle=\ttfamily\small,frame=single]
// Find the node *after which* `value` should go,
// or nullptr if it belongs at the head.
Node* findInsertPos(List& L, int value) {
    Node* before = nullptr;
    Node* cur    = L.head;
    while (cur && value > cur->data) {
        before = cur;
        cur    = cur->next;
    }
    return before;
}
\end{lstlisting}
\end{examplebox}

\subsection*{Cost}

\begin{center}
\begin{tabular}{|l|c|c|l|}
\hline
\textbf{Variant} & \textbf{Best} & \textbf{Average/Worst} & \textbf{Why} \\
\hline
DLL insertion sort & $O(n)$    & $O(n^2)$ & Already-sorted skips inner loop \\
SLL insertion sort & $O(n)$    & $O(n^2)$ & Descending list = trivial \\
\hline
\end{tabular}
\end{center}

The best-case $O(n)$ trigger differs between the two: a sorted DLL
never enters the inner loop (the \texttt{prev} data is already
smaller), while a reverse-sorted SLL inserts every element at the
head --- \texttt{findInsertPos} returns \texttt{nullptr}
immediately, which is also a no-op.

\subsection*{Which array sorts survive the linked-list trip?}

\begin{center}
\begin{tabular}{|l|c|l|}
\hline
\textbf{Algorithm} & \textbf{Port?} & \textbf{Reason} \\
\hline
Insertion sort & Yes      & Only needs neighbor access \\
Merge sort     & Yes      & Splitting \& merging is pointer surgery \\
Selection sort & Yes (ish) & Needs $n$ linear ``find min'' passes \\
Shell sort     & No       & Gap-jumps require indexed access \\
Quicksort      & Poor     & Partition needs both-ends access \\
Heap sort      & No       & Heap needs parent/child by index \\
\hline
\end{tabular}
\end{center}

\begin{ideabox}[Why merge sort is the linked list sort]
Merge sort is the textbook recommendation for linked lists.
Splitting the list requires one walk to find the middle
(\texttt{O(n)}, fine). Merging two sorted lists is pure pointer
splicing: no array index math, no auxiliary buffer, no random
access. Total cost: $O(n \log n)$ time, $O(\log n)$ stack, $O(1)$
extra heap. That is the \emph{same} asymptotic as array merge
sort \emph{without} the $O(n)$ scratch array. Linked lists don't
win many contests, but they win this one.
\end{ideabox}

\begin{notebox}[What the STL actually does]
\texttt{std::list::sort} is merge sort --- specifically a bottom-
up, iterative merge sort that avoids recursion overhead.
\texttt{std::forward\_list::sort} uses the same approach.
\texttt{std::sort} on a \texttt{std::vector} is introsort (quick
$\to$ heap fallback). The STL is explicit: you get different
algorithms on different containers because the container shape
determines what's cheap.
\end{notebox}

\section{4.11 Sentinel (Dummy) Nodes}

Every removal and insertion algorithm you've written so far has a
prologue: a cluster of \texttt{if} statements for the empty list,
the head, and the tail. Sentinel nodes --- also called ``dummy''
nodes --- are a structural trick that makes those special cases
\emph{disappear}. You pay one pointer of overhead and a pinch of
conceptual weirdness; in return, every other method gets shorter.

\begin{defnbox}[Sentinel / dummy node]
A permanent, data-less node that lives at the head (and optionally
the tail) of a list. It is never returned to callers, never counts
toward the list's size, and is never removed. Its purpose is
purely structural: to guarantee that the ``real'' nodes always
have a neighbor on one or both sides.
\end{defnbox}

\subsection*{Why the head-and-tail cases exist}

Take a close look at \texttt{ListRemove} from section 4.8: the
very first lines check \texttt{if (curNode == list->head)} and
\texttt{if (curNode == list->tail)}. Those branches exist because
\emph{head has no predecessor} and \emph{tail has no successor}.
Most of the body is just ``update the list handle instead of an
in-list pointer.''

A sentinel makes head-has-no-predecessor false. Every real node
now has a valid \texttt{prev} pointing to either another real
node or the sentinel. The removal code never has to ask ``am I at
the head?'' because it always has a predecessor to splice around.

\begin{ideabox}[The pattern]
Special cases at the boundary usually mean the data structure is
missing something at the boundary. Sentinels are the canonical
fix: add a virtual neighbor so the boundary stops being a
boundary. This same idea shows up again in trees (nil/leaf
sentinels in red-black trees), in graph algorithms (super-source
/ super-sink nodes), and in strings (null terminators). Once you
see it once you start seeing it everywhere.
\end{ideabox}

\subsection*{One dummy vs. two dummies}

\begin{center}
\begin{tabular}{|l|l|l|}
\hline
\textbf{Variant} & \textbf{Simplifies} & \textbf{Still special-cases} \\
\hline
No sentinel        & ---          & head, tail, empty \\
One head sentinel  & head/empty   & tail \\
Two sentinels      & everything   & none (for interior ops) \\
\hline
\end{tabular}
\end{center}

The two-sentinel DLL is the cleanest: \texttt{ListRemove} is just
``unsplice,'' with no branches at all, because both neighbors are
guaranteed to exist.

\begin{examplebox}[Two-sentinel DLL remove --- branchless]
\begin{lstlisting}[language=C++,basicstyle=\ttfamily\small,frame=single]
void remove(DNode* cur) {
    // cur is a real node, so cur->prev and cur->next
    // are both non-null (they point at real nodes
    // or at the sentinels).
    cur->prev->next = cur->next;
    cur->next->prev = cur->prev;
    delete cur;
}
\end{lstlisting}
\end{examplebox}

Compare that with the no-sentinel version's four pointer checks
(section 4.8). Both are $O(1)$ --- but the branchless version is
faster in practice (the CPU loves straight-line code), shorter,
and almost impossible to get wrong.

\subsection*{What sentinels cost}

\begin{center}
\begin{tabular}{|l|c|l|}
\hline
\textbf{Cost} & \textbf{Amount} & \textbf{Notes} \\
\hline
Extra memory     & 1 or 2 nodes & Constant, regardless of list size \\
Extra complexity & ``size'' is tracked separately & \texttt{size() != number of nodes} \\
Extra care       & iterators must skip sentinels & Usually hidden in \texttt{begin()}/\texttt{end()} \\
\hline
\end{tabular}
\end{center}

The ``size is tracked separately'' point matters for correctness.
A sentinel-based list must either keep a size counter or define
``size'' as ``number of nodes minus the sentinels.'' Walking the
whole list and counting no longer gives the right answer for a
caller who expects 0 on an empty list.

\begin{warnbox}[The empty-list invariant flips]
Without sentinels, an empty list is \texttt{head == nullptr}. With
two sentinels, an empty list has \texttt{head->next == tail} (or
equivalently \texttt{tail->prev == head}). If you inherit sentinel
code without noticing, an empty-list check like
\texttt{if (L.head == nullptr)} silently becomes ``is this list
broken?'' rather than ``is this list empty?'' Update every such
check to the new invariant when you add sentinels.
\end{warnbox}

\subsection*{Do real C++ lists use sentinels?}

Yes. \texttt{std::list} in libstdc++ and libc++ is a doubly-linked
circular list with a single sentinel. The list object itself
\emph{contains} the sentinel inline (no heap allocation for
an empty list), and \texttt{end()} returns an iterator to that
sentinel. That's why dereferencing \texttt{end()} is undefined
behavior but comparing against it is fine: the sentinel exists,
it just doesn't hold a user value.

\begin{notebox}[Circular + sentinel = extremely clean]
libstdc++'s \texttt{std::list} is circular: \texttt{tail->next} is
the sentinel, \texttt{sentinel->next} is the head, and
\texttt{head->prev} is the sentinel. With one sentinel you get
two-sentinel semantics for free, because the sentinel is
simultaneously ``before head'' and ``after tail.'' If you're
curious about the STL source, this is a good place to look ---
it's short, clever, and illustrative.
\end{notebox}

\begin{ideabox}[When to use sentinels in your own code]
For one-off assignment code: probably don't bother --- the
branches you save aren't worth the conceptual weight. For
library-quality data structures: almost always yes. The
simplification to \emph{every} algorithm (insert, remove,
traverse, even iterators) more than pays back the one-time
complexity of ``there's a fake node at each end.''
\end{ideabox}

\subsection*{The payoff: \texttt{std::list::splice}}

The circular-sentinel layout enables one operation that no
contiguous container can match: moving an entire range of nodes
between two lists in $O(1)$, regardless of range size. The STL
exposes this as \texttt{std::list::splice}.

\begin{examplebox}[\texttt{std::list::splice}: the operation only a linked list can do in $O(1)$]
\begin{lstlisting}[language=C++,basicstyle=\ttfamily\small,frame=single]
#include <list>
std::list<int> a {1, 2, 3, 4, 5};
std::list<int> b {10, 20, 30};

// Move the [2, 3, 4] range from a into b before
// the position of 20. O(1) -- no copies, no allocation.
auto first = std::next(a.begin());      // points to 2
auto last  = std::next(a.begin(), 4);   // points to 5 (one past 4)
b.splice(std::next(b.begin()), a, first, last);

// a is now {1, 5}; b is now {10, 2, 3, 4, 20, 30}.
\end{lstlisting}
\end{examplebox}

The cost is $O(1)$ because splice only rewires four pointers ---
the next/prev of the boundary nodes on each side. The range itself
isn't touched. \texttt{std::vector} cannot do this without copying
or moving every element in the range, an inherent $O(k)$ cost. This
is the canonical answer to \emph{``why would I ever use
\texttt{std::list} when vector usually wins?''} If your workload
splices large ranges between lists, the linked structure pays for
its cache-hostility and per-node allocation overhead. Otherwise it
usually doesn't.

\section{4.12 Recursion on Linked Lists}

A linked list is a recursive data structure in disguise: it is
``a head node plus the rest of the list, which is also a linked
list.'' That definition practically hands you the recursive
algorithms for free. Every iterative walker from the last few
sections has a recursive twin that is often shorter, and
\emph{occasionally} clearer --- at a real performance cost.

\begin{ideabox}[The recursive mental model]
Don't think ``the function calls itself.'' Think ``the function
assumes its job on a smaller list is already done, and only has
to handle one more node.'' If the list is empty, the work is
zero. Otherwise, do the head, then trust the recursive call to
handle the tail. That mindset makes the code write itself.
\end{ideabox}

\subsection*{Forward traversal, recursive}

\begin{examplebox}[Recursive forward traverse]
\begin{lstlisting}[language=C++,basicstyle=\ttfamily\small,frame=single]
void traverse(Node* node) {
    if (!node) return;       // empty list: done
    visit(node);             // head: do the work
    traverse(node->next);    // tail: trust recursion
}
\end{lstlisting}
\end{examplebox}

Three lines, one branch. The base case is ``empty list,'' which
maps exactly to \texttt{node == nullptr}. The inductive step is
``visit head, recurse on tail.'' This is the cleanest expression
of a traversal you can write --- and the iterative version from
section 4.9 is still faster in practice.

\subsection*{Forward vs. reverse: one line swap}

The striking thing about recursive list code is how little has to
change to reverse the order. Swap ``visit'' and ``recurse'' and
you get a reverse traversal:

\begin{examplebox}[Recursive reverse traverse]
\begin{lstlisting}[language=C++,basicstyle=\ttfamily\small,frame=single]
void traverseReverse(Node* node) {
    if (!node) return;
    traverseReverse(node->next);  // recurse FIRST
    visit(node);                  // then visit
}
\end{lstlisting}
\end{examplebox}

This works on a \emph{singly}-linked list --- no \texttt{prev}
pointer needed. The reversal is hiding in the call stack: each
activation records a node, all activations unwind in LIFO order,
and visits happen on the way out.

\begin{ideabox}[Recursion \emph{is} a stack]
The call stack is a stack. Recursion that does work ``after the
recursive call'' is literally pushing things onto a stack and
processing them in reverse. This is why reverse traversal on an
SLL ``costs nothing extra'' recursively but $O(n)$ space
iteratively --- the cost didn't vanish, it's just hidden in the
call stack. Section 4.14 will make this explicit by building a
stack ADT.
\end{ideabox}

\subsection*{Recursive search}

\begin{examplebox}[Recursive search]
\begin{lstlisting}[language=C++,basicstyle=\ttfamily\small,frame=single]
Node* search(Node* node, int key) {
    if (!node)              return nullptr; // not found
    if (node->data == key)  return node;    // found
    return search(node->next, key);         // keep looking
}
\end{lstlisting}
\end{examplebox}

Three cases, each a single line. The third case is a \emph{tail
call}: the recursive call is the last thing the function does.
Some compilers (including modern g++/clang with \texttt{-O2}) can
recognize this and rewrite it as a loop, reclaiming the stack
frame each step. With tail-call optimization, the recursive
version has the same stack cost as the iterative one. Without
it, a 10,000-element list will blow the default 8MB stack on
Linux somewhere around recursion depth 100,000 --- well above
10k, but closer than you think.

\subsection*{Cost: recursive vs iterative}

\begin{center}
\begin{tabular}{|l|c|c|l|}
\hline
\textbf{Operation} & \textbf{Rec. time} & \textbf{Rec. space} & \textbf{Iterative space} \\
\hline
Forward traverse  & $\Theta(n)$ & $\Theta(n)^*$ & $\Theta(1)$ \\
Reverse traverse  & $\Theta(n)$ & $\Theta(n)$   & $\Theta(n)$ (explicit stack) \\
Search            & $\Theta(n)$ & $\Theta(n)^*$ & $\Theta(1)$ \\
\hline
\end{tabular}
\end{center}

\noindent$^*$ With tail-call optimization, $\Theta(1)$. C++ does
\emph{not} guarantee TCO --- treat it as a performance
optimization, not a correctness property.

\begin{warnbox}[Stack overflow is a real risk]
A long linked list plus a recursive traversal plus a compiler
that didn't eliminate the tail call = stack overflow. The default
thread stack on Linux is 8MB; on Windows, 1MB; on embedded
systems, often kilobytes. A stack overflow is usually a hard
crash with no useful error message. For production code on
arbitrarily-sized input, prefer iteration. For algorithms where
the recursion depth is bounded (like tree traversals with
$O(\log n)$ depth), recursion is fine.
\end{warnbox}

\begin{notebox}[When recursion is genuinely better]
Flat lists rarely benefit from recursion --- iteration wins on
every axis except elegance. Where recursion pays off: trees
(chapter 6), graphs, divide-and-conquer sorts (merge sort's
``split the list, sort each half''). The unifying property is
branching: when each step produces multiple subproblems, the
recursive structure is essential. One-way linked lists branch
only once per step, which is why iteration is almost always the
right tool here.
\end{notebox}

\begin{ideabox}[Before you move on]
Write both versions once for each algorithm in this chapter.
Not because you'll use the recursive ones in real code, but
because the back-and-forth forces you to see the same algorithm
from two angles. The loop variant and the recursive variant are
always the same algorithm --- only the bookkeeping differs.
Once you can translate between them freely, the rest of the
course (especially trees and graph traversal) becomes much
easier.
\end{ideabox}

\section{4.13 Stack ADT}

Stacks are the first ``restricted list'' ADT. Where a list lets
you touch any element, a stack lets you touch \emph{one}: the
top. That single restriction is the entire point. It buys you a
mental model (LIFO, last-in-first-out), a suite of cheap
operations (everything $O(1)$), and a surprisingly long list of
real-world algorithms that just fall out of it.

\begin{defnbox}[Stack]
A linear ADT with access restricted to one end, called the
\emph{top}. Insertion (\emph{push}) and removal (\emph{pop})
happen only at the top, so the most recently inserted element is
always the first to come out --- last-in, first-out (LIFO).
\end{defnbox}

\subsection*{The five operations}

\begin{center}
\begin{tabular}{|l|l|c|}
\hline
\textbf{Op} & \textbf{Semantics} & \textbf{Cost} \\
\hline
\texttt{push(x)}  & Add \texttt{x} at top           & $O(1)$ \\
\texttt{pop()}    & Remove and return top           & $O(1)$ \\
\texttt{peek()}   & Return top without removing     & $O(1)$ \\
\texttt{empty()}  & True iff no elements            & $O(1)$ \\
\texttt{size()}   & Number of elements              & $O(1)$ \\
\hline
\end{tabular}
\end{center}

All five are \emph{guaranteed} $O(1)$ --- that's not a coincidence,
it's the contract. A data structure that calls itself a stack but
has $O(n)$ push is mislabeled; the whole reason stacks exist is to
be cheap.

\begin{warnbox}[Pop/peek on empty is undefined]
The ADT contract says nothing about what \texttt{pop()} or
\texttt{peek()} does on an empty stack. Different implementations
choose differently: \texttt{std::stack::pop} is UB (nothing
returned, no throw); Python's \texttt{list.pop()} raises
\texttt{IndexError}; some textbook versions return a sentinel.
Always check \texttt{empty()} first, or use whatever
error-signaling mechanism your library provides. Treat this like
dereferencing a null pointer --- the language doesn't stop you,
but the consequences are yours.
\end{warnbox}

\subsection*{C++: \texttt{std::stack}}

C++ provides \texttt{std::stack} as a container \emph{adapter} ---
a wrapper around an underlying container that exposes only the
stack operations. The underlying container defaults to
\texttt{std::deque} but can be any container that supports
\texttt{push\_back}, \texttt{pop\_back}, and \texttt{back}.

\begin{examplebox}[\texttt{std::stack} basics]
\begin{lstlisting}[language=C++,basicstyle=\ttfamily\small,frame=single]
#include <stack>
#include <iostream>

std::stack<int> s;
s.push(7); s.push(14); s.push(9); s.push(5);

while (!s.empty()) {
    std::cout << s.top() << ' ';  // peek
    s.pop();                      // remove
}
// Output: 5 9 14 7   (LIFO)
\end{lstlisting}
\end{examplebox}

Notice \texttt{top()} is the STL's name for \texttt{peek()}, and
\texttt{pop()} returns \texttt{void} --- you must call
\texttt{top()} before \texttt{pop()} if you want the value. This
split is deliberate: it avoids forcing a copy on every pop, which
matters for stacks of expensive-to-copy types.

\begin{ideabox}[When do you pick stack over list?]
Use a stack when your algorithm only ever needs the most-recent
item. That sounds restrictive, but it's \emph{exactly} the shape
of: function calls (the call stack), parsing (parenthesis
matching, shunting-yard), undo/redo, backtracking search, DFS
graph traversal, expression evaluation, and reversing a
sequence. If you catch yourself using \texttt{push\_back} /
\texttt{pop\_back} on a \texttt{vector} and never touching any
middle element, you have a stack --- consider using
\texttt{std::stack} to make the intent obvious and prevent
accidental misuse.
\end{ideabox}

\subsection*{Why restrict access at all?}

A stack has strictly \emph{fewer} operations than a list. That
seems strictly worse. The payoff is threefold:

\begin{enumerate}[leftmargin=*]
    \item \textbf{Speed guarantee.} A list has $O(n)$ random access;
          a stack's operations are all $O(1)$. The restriction is
          what makes that guarantee possible.
    \item \textbf{Code clarity.} A function that takes a
          \texttt{stack<T>\&} advertises its intent: ``I will use
          this LIFO.'' A function taking \texttt{vector<T>\&} is
          making no such promise.
    \item \textbf{Implementation freedom.} Once you commit to LIFO
          semantics, the implementer can swap in any underlying
          container (array, DLL, circular buffer) without
          breaking callers.
\end{enumerate}

\begin{notebox}[The stack is the call stack is the stack]
The term ``stack'' in ``call stack'' isn't a coincidence --- it
\emph{is} the same ADT. When a function calls another, a stack
frame is pushed. When it returns, the frame is popped. The
recursive reverse-traversal from section 4.12 worked because the
call stack \emph{is} a stack, and stacks reverse things. Any
iterative algorithm you can write with an explicit stack, you can
also write recursively (and vice versa). This equivalence is the
backbone of most tree and graph algorithms in later chapters.
\end{notebox}

\section{4.14 Stacks via Linked Lists}

Section 4.13 defined what a stack is. This section picks one
concrete implementation: a singly-linked list where the head
\emph{is} the top. Every stack operation becomes a list operation
you already know.

\subsection*{The mapping}

\begin{center}
\begin{tabular}{|l|l|c|}
\hline
\textbf{Stack op} & \textbf{List op (head = top)} & \textbf{Cost} \\
\hline
\texttt{push(x)}  & \texttt{prepend(x)}           & $O(1)$ \\
\texttt{pop()}    & Read \texttt{head}, remove head & $O(1)$ \\
\texttt{peek()}   & Read \texttt{head->data}      & $O(1)$ \\
\texttt{empty()}  & \texttt{head == nullptr}      & $O(1)$ \\
\texttt{size()}   & Counter updated on push/pop   & $O(1)$ \\
\hline
\end{tabular}
\end{center}

The key design choice is \textbf{head-as-top}. Why not tail?
Because in a \emph{singly}-linked list, prepending is $O(1)$ and
appending-then-removing is $O(n)$ without a tail pointer (you'd
still need to walk to find the new tail after popping). Putting
the top at the head lines the cheap operations up perfectly.

\begin{ideabox}[Choose the end that makes both sides cheap]
For any restricted ADT (stack, queue, deque), the ``which end''
decision is driven by which operations have to stay $O(1)$. For
a singly-linked stack: head, because prepend and remove-head are
both $O(1)$. For a singly-linked queue: head for dequeue,
\emph{tail} for enqueue --- which is why a queue needs a tail
pointer on an SLL but a stack doesn't. Section 4.15 will hit this
exact wall.
\end{ideabox}

\subsection*{Implementation}

\begin{examplebox}[Singly-linked stack]
\begin{lstlisting}[language=C++,basicstyle=\ttfamily\small,frame=single]
template <typename T>
class Stack {
    struct Node { T data; Node* next; };
    Node*  head_ = nullptr;
    size_t n_    = 0;

public:
    ~Stack() { while (!empty()) pop(); }

    bool   empty() const { return head_ == nullptr; }
    size_t size()  const { return n_; }

    void push(const T& x) {
        head_ = new Node{x, head_};  // prepend
        ++n_;
    }

    T pop() {
        // caller's responsibility: !empty()
        Node*  top = head_;
        T      val = std::move(top->data);
        head_      = head_->next;
        delete top;
        --n_;
        return val;
    }

    const T& peek() const { return head_->data; }
};
\end{lstlisting}
\end{examplebox}

A few things worth pointing out:

\begin{itemize}[leftmargin=*]
    \item \textbf{Destructor walks and frees.} A stack owns its
          nodes. Without the destructor, every \texttt{Stack}
          that goes out of scope while non-empty leaks the whole
          chain.
    \item \textbf{\texttt{n\_} counter.} Keeping \texttt{size()}
          $O(1)$ requires maintaining a counter. Walking the list
          every time would make \texttt{size()} $O(n)$, violating
          the ADT contract.
    \item \textbf{\texttt{pop} returns by value; \texttt{peek} by
          const-reference.} Pop has to return because the node is
          about to be deleted, so you can't hand out a reference
          into freed memory. Peek can return a reference because
          the node is still alive.
\end{itemize}

\begin{warnbox}[Copy/move semantics are easy to get wrong]
A class that owns heap memory via a raw pointer needs the Rule of
Three (destructor, copy constructor, copy assignment) or Rule of
Five (add move versions). If you don't write them, C++ generates
shallow copies --- two stacks will share node pointers and
double-free on destruction. The clean fix: hold nodes with
\texttt{std::unique\_ptr<Node>} so ownership is automatic, or
just delegate to \texttt{std::list} / \texttt{std::forward\_list}
underneath.
\end{warnbox}

\subsection*{Linked list vs. array backing}

A stack doesn't \emph{have} to be a linked list. A dynamic array
(\texttt{std::vector}) with push-back-as-push and pop-back-as-pop
is often faster in practice because the elements are contiguous
and the CPU prefetches. The tradeoff:

\begin{center}
\begin{tabular}{|l|c|c|}
\hline
\textbf{Metric} & \textbf{Linked list} & \textbf{Dynamic array} \\
\hline
\texttt{push} worst case  & $O(1)$        & $O(n)$ (on resize) \\
\texttt{push} amortized   & $O(1)$        & $O(1)$ \\
\texttt{pop}              & $O(1)$        & $O(1)$ \\
Cache friendliness        & Poor          & Excellent \\
Memory overhead / elem    & 1 pointer     & 0--1 slot (slack) \\
Handles huge sizes        & Any           & Any \\
\hline
\end{tabular}
\end{center}

For almost every use case, array-backed is the right default ---
which is why \texttt{std::stack} defaults to \texttt{std::deque}
(array-ish under the hood) rather than \texttt{std::list}. The
linked-list version shines when you need \emph{guaranteed}
worst-case $O(1)$ with no amortization.

\begin{notebox}[Amortization, briefly]
A \texttt{vector} push is $O(1)$ on average even though
occasional resizes are $O(n)$. The trick: when the vector doubles
its capacity, the $n$ work is spread across the next $n$ cheap
pushes, so the \emph{amortized} cost is still $O(1)$. Linked
lists don't need this argument --- every push is $O(1)$, full
stop. For real-time systems where a single $O(n)$ spike would
miss a deadline, the linked list's predictable timing wins.
\end{notebox}

\begin{ideabox}[The bigger lesson]
Notice what just happened: you got stack $\to$ done, for free,
by renaming a handful of list operations. That's the point of
layered ADTs. You design the list once, stress-test it, and then
every restricted list (stack, queue, deque) is a ten-line wrapper
that inherits all the correctness work. Building bigger systems
from smaller, already-correct ones is the core abstraction skill
of this course.
\end{ideabox}

\section{4.15 Queue ADT}

A queue is the stack's twin. Where a stack is LIFO (last-in,
first-out), a queue is \emph{FIFO} (first-in, first-out). Both
are restricted list ADTs with exactly two access points --- the
only difference is whether you remove from the same end you insert
at (stack) or the opposite end (queue).

\begin{defnbox}[Queue]
A linear ADT with insertion at one end (the \emph{rear} /
\emph{back}) and removal at the other (the \emph{front}). Order
is preserved: the item pushed first is the item popped first ---
first-in, first-out (FIFO). British English calls a line at a
shop a ``queue''; that is exactly the mental model.
\end{defnbox}

\subsection*{The five operations}

\begin{center}
\begin{tabular}{|l|l|c|}
\hline
\textbf{Op} & \textbf{Semantics} & \textbf{Cost} \\
\hline
\texttt{push(x)} / \texttt{enqueue} & Add \texttt{x} at rear & $O(1)$ \\
\texttt{pop()}   / \texttt{dequeue} & Remove and return front & $O(1)$ \\
\texttt{peek()}  / \texttt{front}   & Return front without removing & $O(1)$ \\
\texttt{empty()} & True iff no elements & $O(1)$ \\
\texttt{size()}  & Number of elements   & $O(1)$ \\
\hline
\end{tabular}
\end{center}

Every operation is $O(1)$, same as a stack. The restriction
(only two access points) is what makes this achievable --- a
general list cannot guarantee $O(1)$ on both ends without extra
structure.

\begin{ideabox}[Stack vs. queue: a one-line summary]
A stack pops what it just pushed. A queue pops what it pushed
the longest time ago. Everything else about the two ADTs is the
same: same shape (linear), same operation count (two active
ends), same cost ($O(1)$ across the board). The semantic
difference --- LIFO vs. FIFO --- is the only thing that actually
changes.
\end{ideabox}

\subsection*{C++: \texttt{std::queue}}

\texttt{std::queue} mirrors \texttt{std::stack}: a container
adapter with a default underlying container (\texttt{std::deque})
and a minimal interface. The names are slightly different because
a queue has two ends.

\begin{examplebox}[\texttt{std::queue} basics]
\begin{lstlisting}[language=C++,basicstyle=\ttfamily\small,frame=single]
#include <queue>
#include <iostream>

std::queue<int> q;
q.push(7); q.push(14); q.push(9); q.push(5);

while (!q.empty()) {
    std::cout << q.front() << ' ';  // peek
    q.pop();                        // remove
}
// Output: 7 14 9 5   (FIFO --- insertion order)
\end{lstlisting}
\end{examplebox}

Notice the output is in \emph{insertion} order, not reversed like
a stack. Also notice \texttt{std::queue} exposes both
\texttt{front()} and \texttt{back()}; you can peek at either end
but only pop from the front and push to the back.

\subsection*{Where queues show up}

\begin{center}
\begin{tabular}{|l|l|}
\hline
\textbf{Use case} & \textbf{What the queue represents} \\
\hline
BFS graph traversal   & Frontier of nodes to visit \\
Scheduling (OS, jobs) & Tasks waiting to run \\
Producer/consumer     & Work items, producer pushes, consumer pops \\
Message passing       & Inbox (network, IPC) \\
Printer / print spool & Jobs queued for the device \\
Buffered I/O          & Bytes waiting to be read or written \\
\hline
\end{tabular}
\end{center}

The unifying theme: \emph{``things arrive over time; serve them
in arrival order.''} That is a much more common pattern than LIFO,
which is why queues are everywhere the moment you leave single-
threaded algorithm land.

\begin{notebox}[BFS uses a queue; DFS uses a stack]
Graph breadth-first search with a queue explores nearer nodes
first. Depth-first search with a stack (or recursion, which
\emph{is} a stack) explores a path as far as possible before
backtracking. The only difference between the two algorithms is
which ADT holds the frontier. Same code, swap the data structure,
get a different traversal order. This is one of the cleanest
illustrations of why ``choice of ADT'' is an algorithmic choice.
\end{notebox}

\subsection*{Why a queue needs two pointers on a linked list}

On a singly-linked list, pop-from-front is cheap (just advance
\texttt{head}) but push-to-back is $O(n)$ --- unless you keep a
tail pointer. A stack's linked implementation needs only the
head, but a queue needs both head and tail because the two
operations live at opposite ends.

\begin{warnbox}[Don't forget to update tail on the last pop]
When a queue empties, \texttt{head} becomes \texttt{nullptr} but
you must also set \texttt{tail} to \texttt{nullptr} --- otherwise
the next push will dereference a pointer to a freed node. This
is the most common bug in handwritten queue code, and it only
shows up in the empty-then-fill again path, which is easy to
miss in testing. Section 4.16 walks through the full
implementation.
\end{warnbox}

\begin{ideabox}[Preview: the other queues]
The basic FIFO queue is the starting point, not the end. Later
sections and courses will add: \emph{priority queue} (pop the
\emph{smallest} element, not the oldest --- implemented with a
heap), \emph{deque} (push and pop from both ends --- implemented
with a doubly-linked list or a ring buffer), and \emph{circular
queue} (fixed-capacity FIFO using modular arithmetic on an
array). Each is a small variation on the same ``ordered work
list'' theme.
\end{ideabox}

\section{4.16 Queues via Linked Lists}

The queue-via-linked-list implementation is structurally almost
identical to the stack version from 4.14, but with one critical
addition: the queue needs a \emph{tail} pointer so push-to-back
stays $O(1)$.

\subsection*{The mapping}

\begin{center}
\begin{tabular}{|l|l|c|}
\hline
\textbf{Queue op} & \textbf{List op (head = front, tail = rear)} & \textbf{Cost} \\
\hline
\texttt{push(x)}  & \texttt{append(x)} via tail pointer & $O(1)$ \\
\texttt{pop()}    & Read \texttt{head}, remove head    & $O(1)$ \\
\texttt{peek()}   & Read \texttt{head->data}           & $O(1)$ \\
\texttt{empty()}  & \texttt{head == nullptr}           & $O(1)$ \\
\texttt{size()}   & Counter updated on push/pop        & $O(1)$ \\
\hline
\end{tabular}
\end{center}

The design is \textbf{head-as-front, tail-as-rear}. Pop removes
from the head ($O(1)$: advance \texttt{head}). Push adds at the
tail ($O(1)$ \emph{only} if we keep a tail pointer).

\begin{ideabox}[The tail pointer is the whole trick]
Without a tail pointer, push is $O(n)$ --- you'd have to walk the
whole list every time to find the end. \emph{With} a tail
pointer, it's two pointer writes: set the old tail's
\texttt{next} to the new node, then set \texttt{tail} to the new
node. That one extra pointer is the difference between a useful
data structure and a useless one.
\end{ideabox}

\subsection*{Implementation}

\begin{examplebox}[Singly-linked queue with head and tail]
\begin{lstlisting}[language=C++,basicstyle=\ttfamily\small,frame=single]
template <typename T>
class Queue {
    struct Node { T data; Node* next; };
    Node*  head_ = nullptr;
    Node*  tail_ = nullptr;
    size_t n_    = 0;

public:
    ~Queue() { while (!empty()) pop(); }

    bool   empty() const { return head_ == nullptr; }
    size_t size()  const { return n_; }

    void push(const T& x) {
        Node* node = new Node{x, nullptr};
        if (tail_) tail_->next = node;  // link old tail
        else       head_       = node;  // was empty
        tail_ = node;
        ++n_;
    }

    T pop() {
        // caller's responsibility: !empty()
        Node* front = head_;
        T     val   = std::move(front->data);
        head_ = head_->next;
        if (!head_) tail_ = nullptr;    // queue now empty
        delete front;
        --n_;
        return val;
    }

    const T& peek() const { return head_->data; }
};
\end{lstlisting}
\end{examplebox}

Two lines deserve close attention:

\begin{itemize}[leftmargin=*]
    \item \textbf{The empty-push case.} When \texttt{tail\_} is
          \texttt{nullptr}, the queue is empty; the new node is
          simultaneously the head and the tail, so both pointers
          must be set. Missing this turns a push on an empty
          queue into a silent no-op.
    \item \textbf{The empty-pop case.} When the pop removes the
          last element, \texttt{tail\_} must be reset to
          \texttt{nullptr} as well. Missing this leaves a
          dangling pointer: the next push will write to
          \texttt{tail\_->next}, which is a pointer into freed
          memory.
\end{itemize}

\begin{warnbox}[Dangling tail is the classic bug]
The most common bug in handwritten queues is forgetting the
\texttt{if (!head\_) tail\_ = nullptr;} line. The bug does not
trigger on the first few operations. It triggers on ``push, pop,
push'' --- and even then only when the list has reached length
zero. If your test suite never exercises that path you will ship
the bug.
\end{warnbox}

\subsection*{Why not a doubly-linked list?}

A DLL would work, but it's overkill. A queue never traverses
backward, never inserts in the middle, and never removes from
anywhere other than the head. The \texttt{prev} pointer on every
node would be dead weight --- pure memory overhead with no
corresponding operation to justify it.

\begin{center}
\begin{tabular}{|l|l|}
\hline
\textbf{Operation} & \textbf{Needs prev?} \\
\hline
\texttt{push(x)}   & No (append at tail, only need tail->next) \\
\texttt{pop()}     & No (remove head, only need head->next) \\
\texttt{peek()}    & No \\
\hline
\end{tabular}
\end{center}

An SLL with head+tail is the minimum sufficient structure. That's
usually the right way to choose data structures: add capability
up to exactly what the operations demand, no further.

\begin{notebox}[What about array-backed queues?]
\texttt{std::queue}'s default (\texttt{std::deque}) is not a
linked list at all --- it's a segmented array (``deque'' $=$
double-ended queue). For most performance-sensitive work, a
circular buffer (fixed-size array with head/tail indices and
modular arithmetic) beats both the linked list and the deque
because it has zero allocations per push and perfect cache
locality. The tradeoff is a fixed capacity, which linked queues
don't have. Pick based on whether your workload has a known upper
bound on queue size.
\end{notebox}

\subsection*{Implementing the ring buffer}

The notebox above gestures at the ring buffer; it deserves a
concrete implementation, since it is the structure used in OS
kernels for bounded producer-consumer queues, in audio/video
pipelines for fixed-latency buffering, and inside
\texttt{std::deque} block-by-block.

\begin{defnbox}[Ring buffer (circular array queue)]
A fixed-capacity array $Q[0\,..\,n-1]$ with two integer indices,
\texttt{head} and \texttt{tail}, that wrap modulo $n$. Items live
in $Q[\texttt{head}], Q[\texttt{head}+1], \ldots, Q[\texttt{tail}-1]$
in circular order. Enqueue writes at $\texttt{tail}$ and advances
$\texttt{tail} \gets (\texttt{tail}+1) \bmod n$; dequeue returns
$Q[\texttt{head}]$ and advances $\texttt{head} \gets
(\texttt{head}+1) \bmod n$. Both are $O(1)$ with zero allocation.
\end{defnbox}

\begin{examplebox}[Ring-buffer queue in C++]
\begin{lstlisting}[language=C++,basicstyle=\ttfamily\small,frame=single]
template <typename T, size_t N>
class RingQueue {
    T   buf_[N];
    size_t head_ = 0;
    size_t tail_ = 0;
    size_t sz_   = 0;     // explicit size: avoids head==tail ambiguity
public:
    bool empty() const { return sz_ == 0; }
    bool full()  const { return sz_ == N; }
    size_t size() const { return sz_; }

    bool enqueue(const T& x) {
        if (full()) return false;     // caller decides: drop or block
        buf_[tail_] = x;
        tail_ = (tail_ + 1) % N;
        ++sz_;
        return true;
    }

    bool dequeue(T& out) {
        if (empty()) return false;
        out  = buf_[head_];
        head_ = (head_ + 1) % N;
        --sz_;
        return true;
    }
};
\end{lstlisting}
\end{examplebox}

\begin{warnbox}[Full-vs-empty disambiguation]
$\texttt{head} == \texttt{tail}$ is ambiguous: it means \emph{both}
``empty'' and ``full'' (after $n$ enqueues with no dequeues, the
write pointer has wrapped all the way around to meet the read
pointer). Two standard fixes:
\begin{itemize}[leftmargin=*]
  \item \textbf{Track size explicitly} (used above): one extra
        \texttt{size\_t sz\_}, no wasted slots. Simplest.
  \item \textbf{Leave one slot unused}: declare full when
        $(\texttt{tail}+1) \bmod n == \texttt{head}$. Saves the size
        counter but capacity becomes $n-1$. CLRS Ch 10.1 uses this.
\end{itemize}
Pick one; never pick neither.
\end{warnbox}

\begin{notebox}[Why ring buffers dominate in low-latency code]
The ring buffer's appeal is structural, not just performance:
fixed capacity at compile time, contiguous memory (cache-friendly),
zero heap allocation per operation, and a trivially mappable
single-producer-single-consumer lock-free variant (separate cache
lines for head and tail, atomic loads/stores --- no mutex). When
queue size has a known upper bound and latency matters, this is
the implementation to reach for, not the linked-list version of
\textsection 4.16 and not \texttt{std::queue}.
\end{notebox}

\begin{ideabox}[Comparing stack 4.14 to queue 4.16]
Side by side: the stack's push was \texttt{prepend}, its pop was
\texttt{remove-head}. The queue's push is \texttt{append}, its
pop is \texttt{remove-head}. The \emph{only} structural
difference between the two is which end the push happens at ---
everything else (pop location, head pointer, size counter, empty
check, destructor) is identical. The two ADTs are more similar
than they are different; the semantics (LIFO vs. FIFO) emerge
entirely from that one choice of end.
\end{ideabox}

\section{4.17 Deque ADT}

A \emph{deque} (pronounced ``deck,'' short for \emph{d}ouble-
\emph{e}nded \emph{que}ue) is the union of a stack and a queue.
Where a stack restricts you to one end and a queue to two specific
ends (front for pop, back for push), a deque lets you push and pop
at \emph{either} end. That flexibility costs nothing in asymptotic
time: all four operations remain $O(1)$.

\begin{defnbox}[Deque]
A linear ADT supporting insertion and removal at both the front
and the back, each in $O(1)$. A deque generalizes both stacks
(use only one end) and queues (push at one end, pop at the other)
in a single structure.
\end{defnbox}

\subsection*{The eight operations}

\begin{center}
\begin{tabular}{|l|l|c|}
\hline
\textbf{Op} & \textbf{Semantics} & \textbf{Cost} \\
\hline
\texttt{push\_front(x)}  & Insert \texttt{x} at front       & $O(1)$ \\
\texttt{push\_back(x)}   & Insert \texttt{x} at back        & $O(1)$ \\
\texttt{pop\_front()}    & Remove and return front          & $O(1)$ \\
\texttt{pop\_back()}     & Remove and return back           & $O(1)$ \\
\texttt{peek\_front()}   & Return front without removing    & $O(1)$ \\
\texttt{peek\_back()}    & Return back without removing     & $O(1)$ \\
\texttt{empty()}         & True iff no elements             & $O(1)$ \\
\texttt{size()}          & Number of elements               & $O(1)$ \\
\hline
\end{tabular}
\end{center}

\subsection*{The ADT hierarchy, visualized}

\begin{center}
\begin{tabular}{|l|c|c|c|c|}
\hline
\textbf{ADT} & \textbf{Push front} & \textbf{Push back} & \textbf{Pop front} & \textbf{Pop back} \\
\hline
Stack  &            &     $\bullet$     &             &     $\bullet$     \\
Queue  &            &     $\bullet$     &  $\bullet$ &                    \\
Deque  & $\bullet$ &     $\bullet$     &  $\bullet$ &     $\bullet$     \\
\hline
\end{tabular}
\end{center}

A stack is a deque that only uses one end. A queue is a deque that
uses opposite ends. \emph{Every} stack and queue operation is also
a valid deque operation --- deques strictly generalize both.

\begin{ideabox}[Why not just always use a deque?]
Because the restrictions on stacks and queues are \emph{features},
not bugs. A function that takes a \texttt{stack\&} is declaring
``I will only ever push/pop from the top''; the type system
enforces that. Swap in a deque and you've invited bugs where
future maintainers accidentally use \texttt{push\_front} and
break the LIFO invariant. Use the most restrictive ADT that
satisfies your algorithm --- it prevents misuse and
self-documents intent.
\end{ideabox}

\subsection*{C++: \texttt{std::deque}}

\texttt{std::deque} is a full container, not an adapter. It
supports random-access iterators, \texttt{operator[]}, and the
standard container interface on top of the eight deque ops.

\begin{examplebox}[\texttt{std::deque} basics]
\begin{lstlisting}[language=C++,basicstyle=\ttfamily\small,frame=single]
#include <deque>
#include <iostream>

std::deque<int> d;
d.push_back(7);     // d: 7
d.push_front(14);   // d: 14 7
d.push_front(9);    // d: 9 14 7
d.push_back(5);     // d: 9 14 7 5

std::cout << d.front() << ' ' << d.back() << "\n"; // 9 5
d.pop_back();       // d: 9 14 7
d.pop_front();      // d: 14 7
std::cout << d[0] << "\n"; // 14  --- random access works
\end{lstlisting}
\end{examplebox}

\begin{notebox}[\texttt{std::deque} is \emph{not} a linked list]
Despite ``deque'' suggesting a linked-list double-ended queue,
\texttt{std::deque} in libstdc++/libc++ is a segmented array
--- a dynamic array of fixed-size chunks. Front and back push
allocate a chunk at most occasionally (amortized $O(1)$), and
random access is $O(1)$ because the chunk pointer array acts as
a two-level index. This is why \texttt{std::stack} and
\texttt{std::queue} default to \texttt{std::deque}: it gives them
cheap insertion/removal at both ends with better cache behavior
than a linked list.
\end{notebox}

\subsection*{Linked-list implementation: the DLL's payoff}

A deque on a linked list is where the doubly-linked list finally
earns its keep. Think about it:

\begin{center}
\begin{tabular}{|l|c|c|}
\hline
\textbf{Operation} & \textbf{SLL cost} & \textbf{DLL cost} \\
\hline
\texttt{push\_front}  & $O(1)$ & $O(1)$ \\
\texttt{push\_back}   & $O(1)$ (with tail) & $O(1)$ \\
\texttt{pop\_front}   & $O(1)$ & $O(1)$ \\
\texttt{pop\_back}    & $O(n)$ (no prev!)  & $O(1)$ \\
\hline
\end{tabular}
\end{center}

The fourth row is the killer. To pop the back of a singly-linked
list, you have to walk the list to find the new tail (the
second-to-last node), because there's no \texttt{prev} pointer to
back up from the current tail. That walk is $O(n)$ --- game over
for the ADT contract.

\begin{warnbox}[You cannot build an $O(1)$ deque on an SLL]
Not even with a tail pointer. Not even with clever bookkeeping.
The pop-back operation fundamentally requires reaching the
second-to-last node, which on an SLL means walking from the head.
If anyone suggests an ``SLL-based deque,'' they either have an
amortized scheme in mind or they have not thought it through. A
deque demands a DLL (or an array).
\end{warnbox}

\begin{ideabox}[Why this chapter built toward this section]
Look back at the order of topics: list $\to$ SLL $\to$ DLL $\to$
stack $\to$ queue $\to$ deque. That progression isn't arbitrary.
The stack worked on an SLL. The queue worked on an SLL (with a
tail pointer). The deque \emph{requires} a DLL to preserve its
cost contract --- the symmetric access pattern on both ends needs
symmetric pointers on every node. This section is where all the
``extra pointer per node'' discussion from 4.6 finally earns its
cost.
\end{ideabox}

\section{4.18 Array-Based Lists}

The chapter has spent 17 sections on pointer-based lists. This
one finally asks: \emph{what if we implemented the list ADT on a
plain old array instead?} The answer reshuffles the cost table.
Array-based lists win on random access and cache performance, and
lose on insertion/removal at the front. Neither approach is
universally better --- they have complementary strengths.

\begin{defnbox}[Array-based list]
A list ADT implemented on top of a contiguous array. A
\emph{dynamic} array-based list allocates more space than the
current length, and resizes (reallocates larger, copies) when
full. A \texttt{length} counter tracks how many slots are in use;
the gap between \texttt{length} and \texttt{capacity} is the
available headroom.
\end{defnbox}

\subsection*{The layout}

\begin{center}
\begin{tabular}{|c|c|c|c|c|c|c|c|}
\hline
\multicolumn{4}{|c|}{\textbf{used}} & \multicolumn{4}{c|}{\textbf{unused capacity}} \\
\hline
\texttt{[0]} & \texttt{[1]} & \texttt{[2]} & \texttt{[3]} & --- & --- & --- & --- \\
\hline
45 & 84 & 12 & 78 & & & & \\
\hline
\end{tabular}
\end{center}

\noindent \texttt{length = 4}, \texttt{capacity = 8}. Elements
live at indices \texttt{[0..length)}; the rest is reserved space
that doesn't need to be copied on the next append.

\subsection*{Append: amortized $O(1)$}

\begin{examplebox}[Append to a dynamic array list]
\begin{lstlisting}[language=C++,basicstyle=\ttfamily\small,frame=single]
void append(List& L, int value) {
    if (L.length == L.capacity) {
        resize(L, L.capacity == 0 ? 1 : L.capacity * 2);
    }
    L.array[L.length++] = value;
}

void resize(List& L, size_t newCap) {
    int* p = new int[newCap];
    for (size_t i = 0; i < L.length; ++i) p[i] = L.array[i];
    delete[] L.array;
    L.array    = p;
    L.capacity = newCap;
}
\end{lstlisting}
\end{examplebox}

A single append is \emph{usually} $O(1)$: assign into
\texttt{array[length]}, increment length. But when the array is
full, the resize step is $O(n)$: allocate, copy every element,
free the old array. So is append $O(1)$ or $O(n)$?

\begin{ideabox}[Amortized analysis: doubling saves you]
Start with capacity 1. After appending $n$ elements, the total
copy work is $1 + 2 + 4 + \ldots + n/2 + n = 2n - 1$. That is
$O(n)$ work spread across $n$ appends, so the \emph{average} per
append is $O(1)$. This is \textbf{amortized} $O(1)$: any
individual append might be slow, but over a long sequence the
cost averages out. The magic number is the doubling factor ---
grow by a constant factor (not a constant amount) or the
argument breaks.
\end{ideabox}

\begin{warnbox}[Growing by +1 is the classic mistake]
If you resize by adding a fixed amount (say, \texttt{capacity += 10})
every time you fill up, appending $n$ elements costs
$O(n^2)$ total, not $O(n)$. Every tenth append triggers a full
copy, and those copies grow linearly. Always grow by a \emph{factor}
(1.5x, 2x --- both work), never an additive constant.
\end{warnbox}

\subsection*{Prepend and insert: the expensive operations}

Appending is cheap because the free space lives at the end.
Prepending is \textbf{$O(n)$} because every existing element has
to shift right by one to make room at index 0. Inserting after
index $k$ shifts everything in $[k+1, \text{length})$ --- worst
case $O(n)$, best case (insert at the end) $O(1)$.

\begin{examplebox}[Array prepend: linear shift]
\begin{lstlisting}[language=C++,basicstyle=\ttfamily\small,frame=single]
void prepend(List& L, int value) {
    if (L.length == L.capacity) resize(L, L.capacity * 2);
    for (size_t i = L.length; i > 0; --i) {
        L.array[i] = L.array[i - 1];       // O(n) shift
    }
    L.array[0] = value;
    ++L.length;
}
\end{lstlisting}
\end{examplebox}

\subsection*{The full cost comparison}

\begin{center}
\begin{tabular}{|l|c|c|}
\hline
\textbf{Operation} & \textbf{Array-based} & \textbf{Linked list} \\
\hline
\texttt{append}                 & $O(1)^*$ & $O(1)$ (with tail) \\
\texttt{prepend}                & $O(n)$   & $O(1)$ \\
\texttt{insert at index $k$}    & $O(n-k)$ & $O(n)$ to find, $O(1)$ to insert \\
\texttt{remove at index $k$}    & $O(n-k)$ & $O(n)$ to find, $O(1)$ to remove \\
\texttt{access by index}        & $O(1)$   & $O(n)$ \\
\texttt{search by value}        & $O(n)$   & $O(n)$ \\
\texttt{append (at end)}        & $O(1)^*$ & $O(1)$ \\
Cache friendliness              & Excellent & Poor \\
Memory overhead per element     & 0--1 slot  & 1--2 pointers \\
\hline
\end{tabular}
\end{center}

\noindent $^*$ amortized

The two structures have \emph{complementary} strengths. Array-
based wins on index access, cache behavior, and append. Linked
wins on front operations and interior insertion once the position
is known. Picking between them is an algorithmic choice driven
by what operations dominate.

\begin{ideabox}[The default is \texttt{std::vector}]
In practice, C++ programmers reach for \texttt{std::vector}
(dynamic array) first and only use \texttt{std::list} when they
have a specific reason. The reasons: $O(n)$ front/middle
operations are rarely actually the bottleneck, cache performance
matters more than big-O for small-to-medium $n$, and memory
overhead per element matters too. ``Always use a linked list for
frequent insertions'' is a textbook myth that doesn't survive
contact with modern hardware. Profile first.
\end{ideabox}

\begin{notebox}[C++: \texttt{std::vector} is the array-based list]
\texttt{std::vector<T>} is precisely this data structure, with a
few extras: exception safety, move semantics, allocators,
iterators, \texttt{emplace\_back} for in-place construction, and
\texttt{shrink\_to\_fit} for giving memory back. Its growth
factor is implementation-defined --- libstdc++ uses 2x, MSVC
uses 1.5x. The 1.5x choice is marginally more memory-efficient
(the old buffer could be reused after a few rounds) but the
difference rarely matters. For any array-based list task in this
course, use \texttt{std::vector} unless you're explicitly
implementing one from scratch.
\end{notebox}

\begin{warnbox}[Iterator / pointer invalidation on resize]
When a \texttt{std::vector} resizes, every pointer, reference,
and iterator into it becomes invalid --- the old memory is freed
and the elements live at fresh addresses. Code like
\texttt{auto\& x = v[0]; v.push\_back(...); x = 5;} is
undefined behavior if the push triggered a resize. Linked lists
don't have this problem (nodes stay where they were), which is
one real use case for preferring them.
\end{warnbox}

\section*{List-variant decision matrix}

The four realizations of the List ADT covered in this chapter --
\textbf{dynamic array} (\texttt{std::vector}), \textbf{singly linked
list} (SLL), \textbf{doubly linked list} (DLL), and \textbf{sentinel
DLL} (or \texttt{std::list}) -- trade cost in predictable ways. This
matrix is the decision you'll actually make on assignments.

\begin{ideabox}[Cost of each List-ADT operation]
All entries are worst-case unless marked. $n$ = list length.

\vspace{0.4em}
\noindent
\begin{tabular}{@{}l|cccc@{}}
\textbf{Operation}
    & \textbf{Array}
    & \textbf{SLL}
    & \textbf{DLL}
    & \textbf{Sentinel DLL}\\
\hline
Access $v[i]$           & $O(1)$            & $O(n)$            & $O(n)$            & $O(n)$\\
Append (back)           & $O(1)^*$ amort.\  & $O(1)$ w/ tail    & $O(1)$ w/ tail    & $O(1)$\\
Prepend (front)         & $O(n)$ shift      & $O(1)$            & $O(1)$            & $O(1)$\\
Insert after node       & $O(n)$ shift      & $O(1)$            & $O(1)$            & $O(1)$\\
Insert before node      & $O(n)$ shift      & $O(n)$            & $O(1)$            & $O(1)$\\
Remove node (have ptr)  & $O(n)$ shift      & $O(n)$ find prev  & $O(1)$            & $O(1)$\\
Remove by value         & $O(n)$            & $O(n)$            & $O(n)$            & $O(n)$\\
Reverse traversal       & $O(n)$            & $O(n)$ reconstruct& $O(n)$            & $O(n)$\\
Memory per element      & $\sim$ 4--8 B     & + 1 pointer       & + 2 pointers      & + 2 pointers\\
Cache friendliness      & excellent         & poor              & poor              & poor\\
\end{tabular}
\end{ideabox}

\begin{ideabox}[Which variant, for which job?]
Read top-down; take the first match.

\medskip
\textbf{Do you need $O(1)$ indexing ($v[i]$)?}
$\to$ \textbf{array-based} (\texttt{std::vector}). Nothing else gives you
random access. Accept $O(n)$ middle inserts.

\medskip
\textbf{Is the data mostly appended/consumed at the \emph{ends} only?}
(Stack: push/pop one end. Queue: enqueue one end, dequeue the other.
Log buffer: append, scan.)\\
$\to$ \textbf{array-based} for stacks (\texttt{std::stack} $\to$
\texttt{std::deque}). \textbf{Array-based deque} or \textbf{linked list}
for queues. \texttt{std::queue} uses \texttt{std::deque}.

\medskip
\textbf{Are you doing many mid-list inserts/removes, AND you already hold
a pointer/iterator to the node (e.g., LRU cache, intrusive list in
kernel code)?}\\
$\to$ \textbf{DLL} (or \textbf{sentinel DLL}). This is the scenario linked
lists win; every other scenario, arrays usually win.

\medskip
\textbf{Do you need bidirectional traversal?}
$\to$ \textbf{DLL} or \textbf{sentinel DLL}. SLL only walks forward;
walking backward means rebuilding a reverse list.

\medskip
\textbf{Do you want \texttt{insertBefore} as a primitive, or do the
empty/first/last cases in your code feel gross?}\\
$\to$ \textbf{sentinel DLL}. The dummy head/tail nodes make every splice
look the same; no special-casing.

\medskip
\textbf{Is every byte precious (embedded, cold cache)?}
$\to$ \textbf{SLL}, not DLL: one pointer per node, not two. And consider
an \textbf{array-based list} anyway -- contiguous memory often wins even
on ``slow at middle insert'' workloads because of cache.

\medskip
\textbf{Still unsure?}
$\to$ \textbf{array-based list} (\texttt{std::vector}). It's the right
default. Linked lists are special-purpose; reach for them when you have
a concrete reason (pointer stability, intrusive data, O(1) splicing).
\end{ideabox}

\begin{warnbox}[The ``linked lists are always faster for middle insert''
myth]
The Big-O says so; real benchmarks often disagree. Inserting into a
\texttt{std::vector} at position $k$ is $O(n)$ shifts of small contiguous
bytes -- extremely cache-friendly. Inserting into a \texttt{std::list} is
$O(1)$ but requires walking to position $k$ first, pointer-chasing
non-contiguous nodes -- extremely cache-hostile. Below roughly $10^4$
elements, vector insert-at-middle usually wins. Above it, it depends on
access pattern. \textbf{Measure, don't assume.}
\end{warnbox}

\begin{ideabox}[Chapter 4 wrap-up]
You've now seen two complete implementations of the list ADT
--- one pointer-based, one array-based --- and three restricted
variants (stack, queue, deque) layered on top. The unifying
theme: \textbf{cost is a function of layout}, not of semantics.
Two data structures with the same API can have wildly different
performance characteristics because their memory layout makes
different operations cheap. Chapter 5 moves from linear
containers to hierarchical ones (trees), where the layout
decisions get even more interesting.
\end{ideabox}

\section*{What this connects to}

\begin{notebox}[Forward links to later chapters]
\begin{itemize}
    \item \textbf{Chapter 5 (hash tables).} Chaining is implemented as an
    array of \emph{linked lists}: the node + pointer idiom from this
    chapter is reused verbatim as the overflow chain per bucket.
    \item \textbf{Chapter 6 (trees, BSTs).} The node \texttt{struct}
    generalizes from one \texttt{next} pointer to two (\texttt{left},
    \texttt{right}). Tree traversal is linked-list traversal with a
    branching factor of 2.
    \item \textbf{Optional ch.~10 (graphs).} Adjacency-list graph
    representation is \texttt{vector<vector<int>>} -- i.e., a vector of
    lists. The cost arguments here tell you when to prefer adjacency
    matrix instead.
    \item \textbf{Stack / queue reuse everywhere.} DFS uses an explicit
    stack (or the call stack). BFS uses a queue. Expression parsing, undo
    systems, and recursion unwinding all lean on the Stack ADT from 4.13.
\end{itemize}
\end{notebox}

\textbf{Companion materials.} Compact notes:
\texttt{chapters/ch\_4/notes.tex} (includes the decision matrix in compact
form). Practice prompts (problem $\to$ pseudocode $\to$ C++):
\texttt{chapters/ch\_4/practice.md}.

\end{document}
